% ----------------------------------------------------------------
% AMS-LaTeX Paper ************************************************
% **** -----------------------------------------------------------
\documentclass{amsart}
\usepackage [spanish] {babel}
\usepackage{amsmath}
\usepackage[cal=boondox]{mathalfa}
\usepackage{graphicx}
\usepackage{amssymb}
\usepackage{amsbsy} % mas simbolitos
\usepackage{amsfonts}
\usepackage{amsthm}
\usepackage{amsxtra}
\usepackage{dsfont}
\usepackage{upgreek} % para poner letras griegas sin cursiva
\usepackage{cancel} % para tachar
\usepackage{mathdots} % para el comando \iddots
\usepackage{mathrsfs} % para formato de letra
\usepackage{stackrel} % para el comando \stackbin
%\usepackage{doublestroke}
% ----------------------------------------------------------------
\vfuzz2pt % Don't report over-full v-boxes if over-edge is small
\hfuzz2pt % Don't report over-full h-boxes if over-edge is small
% THEOREMS -------------------------------------------------------
\newtheorem{thm}{Theorem}[section]
\newtheorem{cor}[thm]{Corollary}
\newtheorem{lem}[thm]{Lemma}
\newtheorem{prop}[thm]{Proposition}
\theoremstyle{definition}
\newtheorem{defn}[thm]{Definition}
\theoremstyle{remark}
\newtheorem{rem}[thm]{Remark}
\numberwithin{equation}{section}
% MATH -----------------------------------------------------------
\newcommand{\norm}[1]{\left\lVert#1\right\rVert}
\newcommand{\abs}[1]{\left\lvert#1\right\rvert}

\newcommand{\set}[1]{\left\lbrace{#1}\right\rbrace}
\newcommand{\setprop}[2]{\left\lbrace{ #1 }\quad|\quad{ #2 }\right\rbrace}
\newcommand{\setR}{\mathds R}
\newcommand{\setN}{\mathds N}
\newcommand{\setZ}{\mathds Z}
\newcommand{\setQ}{\mathds Q}
\newcommand{\setC}{\mathds C}
\newcommand{\setH}{\mathds H}
\newcommand{\setO}{\mathds O}
\newcommand{\setPrimes}{\mathds P}
\newcommand{\eps}{\varepsilon}
\newcommand{\fromto}{\longrightarrow}
\newcommand{\BX}{\mathbf{B}(X)}
\newcommand{\ssA}{\mathcal{A}}
\newcommand{\setsymbol}[1]{\textsf{#1}}
\newcommand{\sA}{\setsymbol{A}}
\newcommand{\sB}{\setsymbol{B}}
\newcommand{\sC}{\setsymbol{C}}
\newcommand{\sD}{\setsymbol{D}}
\newcommand{\sE}{\setsymbol{E}}
\newcommand{\sF}{\setsymbol{F}}
\newcommand{\sG}{\setsymbol{G}}
\newcommand{\sH}{\setsymbol{H}}
\newcommand{\sI}{\setsymbol{I}}
\newcommand{\sJ}{\setsymbol{J}}
\newcommand{\sK}{\setsymbol{K}}
\newcommand{\sL}{\setsymbol{L}}
\newcommand{\sM}{\setsymbol{M}}
\newcommand{\sN}{\setsymbol{N}}
\newcommand{\sO}{\setsymbol{O}}
\newcommand{\sP}{\setsymbol{P}}
\newcommand{\sQ}{\setsymbol{Q}}
\newcommand{\sR}{\setsymbol{R}}
\newcommand{\sS}{\setsymbol{S}}
\newcommand{\sT}{\setsymbol{T}}
\newcommand{\sU}{\setsymbol{U}}
\newcommand{\sV}{\setsymbol{V}}
\newcommand{\sW}{\setsymbol{W}}
\newcommand{\sX}{\setsymbol{X}}
\newcommand{\sY}{\setsymbol{Y}}
\newcommand{\sZ}{\setsymbol{Z}}
\newcommand{\dsa}{\setsymbol{a}}
\newcommand{\dsb}{\setsymbol{b}}
\newcommand{\dsc}{\setsymbol{c}}
\newcommand{\dsd}{\setsymbol{d}}
\newcommand{\dse}{\setsymbol{e}}
\newcommand{\dsf}{\setsymbol{f}}
\newcommand{\dsg}{\setsymbol{g}}
\newcommand{\dsh}{\setsymbol{h}}
\newcommand{\dsi}{\setsymbol{i}}
\newcommand{\dsj}{\setsymbol{j}}
\newcommand{\dsk}{\setsymbol{k}}
\newcommand{\dsl}{\setsymbol{l}}
\newcommand{\dsm}{\setsymbol{m}}
\newcommand{\dsn}{\setsymbol{n}}
\newcommand{\dso}{\setsymbol{o}}
\newcommand{\dsp}{\setsymbol{p}}
\newcommand{\dsq}{\setsymbol{q}}
\newcommand{\dsr}{\setsymbol{r}}
\newcommand{\dss}{\setsymbol{s}}
\newcommand{\dst}{\setsymbol{t}}
\newcommand{\dsu}{\setsymbol{u}}
\newcommand{\dsv}{\setsymbol{v}}
\newcommand{\dsw}{\setsymbol{w}}
\newcommand{\dsx}{\setsymbol{x}}
\newcommand{\dsy}{\setsymbol{y}}
\newcommand{\dsz}{\setsymbol{z}}
\newcommand{\mthds}[1]{\setsymbol{#1}}

\newcommand{\ssetetsymbol}[1]{\mathcal{#1}}
\newcommand{\ssetA}{\ssetsymbol{A}}
\newcommand{\ssetB}{\ssetsymbol{B}}
\newcommand{\ssetC}{\ssetsymbol{C}}
\newcommand{\ssetD}{\ssetsymbol{D}}
\newcommand{\ssetE}{\ssetsymbol{E}}
\newcommand{\ssetF}{\ssetsymbol{F}}
\newcommand{\ssetG}{\ssetsymbol{G}}
\newcommand{\ssetH}{\ssetsymbol{H}}
\newcommand{\ssetI}{\ssetsymbol{I}}
\newcommand{\ssetJ}{\ssetsymbol{J}}
\newcommand{\ssetK}{\ssetsymbol{K}}
\newcommand{\ssetL}{\ssetsymbol{L}}
\newcommand{\ssetM}{\ssetsymbol{M}}
\newcommand{\ssetN}{\ssetsymbol{N}}
\newcommand{\ssetO}{\ssetsymbol{O}}
\newcommand{\ssetP}{\ssetsymbol{P}}
\newcommand{\ssetQ}{\ssetsymbol{Q}}
\newcommand{\ssetR}{\ssetsymbol{R}}
\newcommand{\ssetS}{\ssetsymbol{S}}
\newcommand{\ssetT}{\ssetsymbol{T}}
\newcommand{\ssetU}{\ssetsymbol{U}}
\newcommand{\ssetV}{\ssetsymbol{V}}
\newcommand{\ssetW}{\ssetsymbol{W}}
\newcommand{\ssetX}{\ssetsymbol{X}}
\newcommand{\ssetY}{\ssetsymbol{Y}}
\newcommand{\ssetZ}{\ssetsymbol{Z}}
\newcommand{\csseta}{\ssetsymbol{a}}
\newcommand{\cssetb}{\ssetsymbol{b}}
\newcommand{\cssetc}{\ssetsymbol{c}}
\newcommand{\cssetd}{\ssetsymbol{d}}
\newcommand{\cssete}{\ssetsymbol{e}}
\newcommand{\cssetf}{\ssetsymbol{f}}
\newcommand{\cssetg}{\ssetsymbol{g}}
\newcommand{\csseth}{\ssetsymbol{h}}
\newcommand{\csseti}{\ssetsymbol{i}}
\newcommand{\cssetj}{\ssetsymbol{j}}
\newcommand{\cssetk}{\ssetsymbol{k}}
\newcommand{\cssetl}{\ssetsymbol{l}}
\newcommand{\cssetm}{\ssetsymbol{m}}
\newcommand{\cssetn}{\ssetsymbol{n}}
\newcommand{\csseto}{\ssetsymbol{o}}
\newcommand{\cssetp}{\ssetsymbol{p}}
\newcommand{\cssetq}{\ssetsymbol{q}}
\newcommand{\cssetr}{\ssetsymbol{r}}
\newcommand{\cssets}{\ssetsymbol{s}}
\newcommand{\cssett}{\ssetsymbol{t}}
\newcommand{\cssetu}{\ssetsymbol{u}}
\newcommand{\cssetv}{\ssetsymbol{v}}
\newcommand{\cssetw}{\ssetsymbol{w}}
\newcommand{\cssetx}{\ssetsymbol{x}}
\newcommand{\cssety}{\ssetsymbol{y}}
\newcommand{\cssetz}{\ssetsymbol{z}}
\newcommand{\mthcl}[1]{\ssetsymbol{#1}}

\newcommand{\card}[1]{\abs{#1}}
\newcommand{\numeral}{\natural}

\newcommand{\monrel}[1]{\mathbf{#1}}
\newcommand{\monrelinv}[1]{\overset{\smallsmile}{\monrel{#1}}}
\newcommand{\monrelat}[2]{\monrel{#1}{\left({#2}\right)}}
\newcommand{\monrelinvat}[2]{\monrelinv{#1}{\left({#2}\right)}}


\newcommand{\binrel}[1]{\monrel{#1}}
\newcommand{\binrelinv}[1]{\overset{\smallsmile}{\binrel{#1}}}
\newcommand{\binrelat}[3]{\binrel{#1}{\left({#2},\;{#3}\right)}}
\newcommand{\binrelinvat}[3]{\binrelinv{#1}{\left({#2},\;{#3}\right)}}


\newcommand{\funct}[1]{\monrel{#1}}
\newcommand{\functinv}[1]{{\funct{#1}}^{-1}}
\newcommand{\functat}[2]{\funct{#1}{\left({#2}\right)}}
\newcommand{\functinvat}[2]{\functinv{#1}{\left({#2}\right)}}


\newcommand{\defin}{\overset {def}{=}}
\newcommand{\leqpot}{\preccurlyeq}
\newcommand{\geqpot}{\succcurlyeq}
\newcommand{\gpot}{\succ}
\newcommand{\lpot}{\prec}
% ----------------------------------------------------------------
\begin{document}
	
	\title{Conjuntos y Números Cardinales}%
	\author{Julián Calderón Almendros}%
	\address{Málaga, España}%
	\email{julian.calderon.almendros@gmail.com}%
	
	\thanks{Alba Calderón Muñoz}%
	\subjclass{Teoría de Conjuntos Intuitiva}%
	\keywords{Conjunto , Cardinal , Número}%
	
	\date{15/07/2023}%
	\dedicatory{A Alba, que me inspiró}%
	%\commby{}%
	% ----------------------------------------------------------------
	
	\begin{abstract}
		Este artículo pretende hacer ver como podemos ver los números de
una manera distinta: a partir de ciertas comparaciones entre colecciones de
objetos. Así, pienso, será más fácil introducir la teoría de conjuntos en
principiantes.
	\end{abstract}
	\maketitle
	% ----------------------------------------------------------------
	\section{Las colecciones (o conjuntos)}
	
	
	Lo primero es darnos cuenta que nuestro entorno más inmediato está
repleto de cosas y de colecciones de cosas. Muchas de las cosas son, de
hecho, colecciones de otras cosas.
	
	
	Pongamos algún ejemplo. Yo estoy ahora mismo escribiendo este artículo.
Delante de mí hay una pantalla, que veo como un todo pero que toma sentido
cuando la vemos como una matriz -colección- de píxeles. Mientras tecleo un
teclado mecánico: tecleo un conjunto ordenado de teclas que se corresponden
con las letras del alfabeto español y algunos signos más. En este documento
hay palabras, un montón. A mi alrededor hay un desorden completo: un montón
de objetos indiferenciados.
	
	
	Las colecciones son objetos habituales de nuestro mundo, o más bien, de
la forma en que observamos y comprendemos el mundo.
	
	
	
	Ahora bien: ¿cómo relacionamos estos conjuntos entre sí? Decimos: -en
esta colección hay más objetos que esta otra colección- o bien -estas dos
colecciones son pares en número de elementos-. Pero aún no hemos introducido
el concepto de número. Solo de colecciones. ¿Cómo podemos establecer una
cierta igualdad entre las colecciones?
	
	
	Llamaremos a dos colecciones, conjuntos, $\sA$ y $\sB$. Que $\sA$ sea un
conjunto quiere decir que tiene ciertos elementos y que podemos al menos
nombrarlos. Veamos como escribirlo. $\sA = \{a , b , c , d , e \}$. Y
escribimos $\sB$ como $\sB = \{ 2 , 4 , 7 , 1 , 0 \}$. Esta será la primera
forma de escribir los conjuntos en detalle. Tiene sus limitaciones (supón que
tienes que describir la colección de palabras distintas que hay en \emph{El
Quijote}), pero para nuestros propósitos valdrá.
	
	
	Algo importante en los conjuntos es que los elementos son únicos.
Pongamos el caso del conjunto $\sS\defin\{1,2,3,4,1,2,3,4\}$ y el conjunto
$\sR\defin\{1,2,3,4\}$ y por último el conjunto $\sP\defin\{4,1,3,2\}$. Pues
son el mismo conjunto. $\sS = \sR = \sP$.
	
	
	¿Porqué son el mismo conjunto? Porque los conjuntos se diferencian solo
en el caso que algún elemento de un conjunto no esté en otro y viceversa.
Pongamos el caso de $\sS$ y $\sR$. 
	
	
	Cojamos elementos de $\sR$. Por ejemplo, por orden de aparición.
$1\in\sR$. Podemos observar que también $1\in\sS$. De acuerdo está dos veces,
pero está.  $2\in\sR$. Podemos observar que también $2\in\sS$. De acuerdo
también está dos veces, pero está en ambos. Así podemos seguir hasta
$4\in\sR$. Podemos observar que también $4\in\sS$. Que también está dos
veces. Para lo que nos importa: está. Como todo elemento de $\sR$ está en
$\sS$ se dice que $\sR\subseteq\sS$ que en palabras se dice que el conjunto
$\sR$ está incluido en $\sS$. Date cuneta que el símbolo elegido para estar
incluido es muy parecido al de menor o igual ($\leq$).
	
	
	Pero, ¿estará el conjunto $\sS$ incluido en $\sR$. Como antes vamos a ir
elemento por elemento de $\sS$. Pero tenemos mucho trabajo hecho: los cuatro
primeros elementos ya sabemos que son compartidos por los dos conjuntos.
Comencemos por el quinto elemento en orden de aparición: $1\in\sS$ que
sabemos que está en $\sR$. Es el primero en $\sR$ en orden de aparición.
Parece que va a ser aburrido. Los siguientes son el 2, el 3 y el 4 que ya
habíamos comprobado anteriormente. Así que efectivamente $\sS$ está incluido
en $\sR$: $\sS\subseteq\sR$.
	
	
	Pero entonces tenemos que $\sR\subseteq\sS$ y $\sS\subseteq\sR$. Todo
elemento de $\sS$ está en $\sR$ y viceversa: todo elemento de $\sR$ está en
$\sS$. No hay diferencias entre ambos conjuntos mirando sus elementos: pero
es que un conjunto solo depende de sus elementos. Concluimos que $\sS=\sR$.
Son exactamente la misma cosa. $\{1,2,3,4\}=\{1,2,3,4,1,2,3,4\}$. Y por más
veces que repitamos un elemento en su representación entre llaves, sigue
siendo exactamente lo mismo.
	
	
	Comparemos ahora los conjuntos $\sR$ y $\sP$. Antes afirmé que se trata
de dos conjuntos idénticos. Y si miramos el razonamiento anterior, que
podemos fácilmente repetir sin cambios, veremos que así es. Primero
llegaremos a que $\sR\subseteq\sP$ y luego a que $\sP\subseteq\sR$.
Finalmente eso equivale a $\sR=\sP$. Y es lógico: tener un elemento en tu
haber no dice nada del puesto en la representación. Nada de nada. O bien el
elemento está o bien el elemento no está, es lo única propiedad que relaciona
un elemento con un conjunto. El elemento $\mathbf{e}$ pertenece al conjunto
$\sE$ o bien no lo hace. Ninguna otra propiedad relaciona a $\mathbf{e}$ y a
$\sE$. De esas dos cosas solo podemos decir que $\mathbf{e}\in\sE$ o
$\mathbf{e}\notin\sE$. Nada más. Un conjunto es el objeto matemático más
sencillo que podemos encontrar.
	
	
	Lo que acabamos de encontrar se dice en símbolos así:
	
	
	Definimos $\sE\subseteq\sF$ como:
	$$\forall \mathbf{e} \;\; \mathbf{e}\in\sE \;\Rightarrow\;
\mathbf{e}\in\sF$$
	$$\text{Para todo elemento }\; \mathbf{e} \;\; \text{ si
}\;\mathbf{e}\in\sE \;\text{entonces}\; \mathbf{e}\in\sF$$
	
	
	Por otra parte hemos llegado a la siguiente (y fundamental) conclusión:
	
	
	Decimos que $\sE = \sF$ cuando:
	$$\forall \mathbf{e} \;\; \mathbf{e}\in\sE \;\Rightarrow\;
\mathbf{e}\in\sF\;\;\;\wedge\;\;\;\forall \mathbf{f} \;\; \mathbf{f}\in\sF
\;\Rightarrow\; \mathbf{f}\in\sE$$
	\begin{gather*}
		\text{Para todo elemento }\; \mathbf{e} \;\; \text{ si
}\;\mathbf{e}\in\sE \;\text{entonces}\; \mathbf{e}\in\sF\\
		\text{ a la vez que } \\
		\text{ para todo elemento }\; \mathbf{f} \;\; \text{ si
}\;\mathbf{f}\in\sF \;\text{entonces}\; \mathbf{f}\in\sE
	\end{gather*}
	
	
	
	Así llegamos a la conclusión que


	$$ \sE = \sF \;\;\;\; \Leftrightarrow \;\;\;\;  \sE\subseteq\sF
\;\;\wedge\;\;\sF\subseteq\sE$$
	\begin{gather*}
		\sE = \sF\\
		\text{ si y solo si } \\
		\sE \text{ está incluido en } \sF \text{ \textbf{y} }\sF \text{
también está incluido en } \sE 
	\end{gather*}
	
	
	
	Advertencia: habrá colecciones (la inmensa mayoría) que no podamos
colocar entre llaves, designando uno por uno sus elementos. Ve pensando
alguna.
	
	
	\section{Las correspondencia entre conjuntos}
	
	
	Cojamos dos conjuntos cualquiera, $\sA$ y $\sB$. Vamos a decir que el
conjunto $\sA$ será nuestro \emph{conjunto de origen en la correspondencia}
(también lo podemos llamar \emph{conjunto de inicial}), correspondencia que
llamaremos $\monrel{f}$, y que $\sB$ será nuestro \emph{conjunto final}. Esto
lo denotaremos así: $\monrel{f} : \sA \fromto \sB$. A veces acortaremos y
diremos solamente $\monrel{f}$, pero esto solo es porque somos un poco vagos
al escribir, el nombre ha de incluir también al conjunto inicial y al final.
	
	
	¿Qué es una correspondencia? Es una relación que comienza en $\sA$ o
\emph{conjunto inicial} y termina en $\sB$ o \emph{conjunto final}. Para
decirlo de una forma coloquial, hay elementos del conjunto inicial que
capturan elementos del conjunto final. Esta captura la expresamos con flechas
con barra vertical en el origen $\mapsto$.
	
	
	
	¿Para qué necesitamos de las correspondencias?
	
	
	Bueno, como veremos dentro de poco, los números se pueden ver como 
relaciones entre colecciones. Las correspondencias hacen el papel de las 
relaciones. Una correspondencia $\monrel{f}:\sA\fromto\sB$ es una relación de 
la colección $\sA$ en la colección $\sB$.  Y ¿cómo es una correspondencia? Ya 
sabemos que los conjuntos son sus elementos y ya está, pues una 
correspondencia entre dos conjuntos será una asociación de elementos del 
conjunto inicial con elementos del conjunto final. Pongamos que la colección 
$\sF=\set{Pepe,Paco,Isabel,Alicia,Alba,Fernando}$ es conjunto inicial y final 
a la vez. Esta será la correspondencia es "ser hermano de" y la denotaremos 
por $\monrel{sh}: \sF\fromto\sF$.
	Concretemos esta correspondencia $\monrel{sh}$:
	
	
	\begin{alignat*}{4}
		\monrel{sh} : &\sF &\fromto &\sF                   \\
		&\text{José}     &\mapsto  &\text{Isabel}          \\
		&\text{José}     &\mapsto  &\text{Alicia}          \\
		&\text{Juan}     &\mapsto  &\text{Alba}            \\
		&\text{Isabel}   &\mapsto  &\text{José}            \\
		&\text{Isabel}   &\mapsto  &\text{Alicia}          \\
		&\text{Alicia}   &\mapsto  &\text{José}            \\
		&\text{Alicia}   &\mapsto  &\text{Isabel}          \\
		&\text{Alba}     &\mapsto  &\text{Juan}            \\
		&\text{Fernando} & \;      &\text{Fernando}
	\end{alignat*}
	
	
	Cualquier combinación de flechas del conjunto inicial llegando al 
conjunto final (o falta de flechas) es válida en una correspondencia. Puede 
no haber ninguna flecha o existir todas las flechas posibles de cada elemento 
del conjunto original a los elementos del conjunto final.
	
	
	Una forma de expresar una correspondencia que se ajusta bastante bien a 
todo lo que acabamos de decir es mediante otro conjunto. Por ejemplo, tomando 
el conjunto anterior, $\sF=\set{Pepe, Paco, Isabel, Alicia, Alba, Fernando}$,
la correspondencia $\monrel{sh}:\sF\fromto\sF$ se puede expresar como
	\begin{alignat*}{10}
		\monrel{sh} : & \sF & \fromto & \sF & :: & \{ && \text{José}   &
 \mapsto & \text{Isabel} & , & \\
		&&&&&&& \text{José}   & \mapsto & \text{Alicia} & ,&\\ 
		&&&&&&& \text{Juan}   & \mapsto & \text{Alba} & ,  &\\
		&&&&&&& \text{Isabel} & \mapsto & \text{José} & ,  &\\
		&&&&&&& \text{Isabel} & \mapsto & \text{Alicia} & ,&\\ 
		&&&&&&& \text{Alicia} & \mapsto & \text{José} & ,  &\\
		&&&&&&& \text{Alicia} & \mapsto & \text{Isabel} & ,&\\
		&&&&&&& \text{Alba}   & \mapsto & \text{Juan}   && \} 
	\end{alignat*} 
	
    
    Como ves, aquí queda reflejada la correspondencia que anteriormente hemos
    pintado completamente, esto es, no falta por expresar nada de la
    correspondencia \emph{ser hermanos}, llamada $\monrel{sh}$ anteriormente
    definida. Ahora voy a transformar un poco la notación: 
	
    
    \begin{alignat*}{13}
		\monrel{sh} : & \sF & \fromto & \sF & :: & \{ && ( & \text{José}   & 
, & \text{Isabel} & ) & , & \\
		&&&&&&& ( & \text{José}   & , & \text{Alicia} & ) & ,&\\ 
		&&&&&&& ( & \text{Juan}   & , & \text{Alba} & ) & ,  &\\
		&&&&&&& ( & \text{Isabel} & , & \text{José} & ) & ,  &\\
		&&&&&&& ( & \text{Isabel} & , & \text{Alicia} & ) & ,&\\ 
		&&&&&&& ( & \text{Alicia} & , & \text{José} & ) & ,  &\\
		&&&&&&& ( & \text{Alicia} & , & \text{Isabel} & ) & ,&\\
		&&&&&&& ( & \text{Alba}   & , & \text{Juan}  & ) & & \} 
	\end{alignat*}  
	
    
    A esta notación los llamaremos por \emph{parejas ordenadas}. El apelativo 
    ordenadas es sencillo de explicar: el primer elemento de la pareja 
    pertenece al conjunto inicial y el segundo al final. Aunque el conjunto 
    inicial y el final sean el mismo sigue teniendo sentido, porque indica la 
    dirección de la flecha. Del primer elemento de la pareja al segundo.
	
	\section{Las aplicaciones entre conjuntos}
	
	No nos vamos a detener en las correspondencias. Hay clases de 
correspondencias muy chulas, esto es, correspondencias que reúnen ciertos 
requisitos que las hacen muy parecidas a cosas como la ordenación o a la 
igualdad. Pero ahora mismo vamos a seguir un camino más llano, más directo: 
nos centraremos en una clase muy especial de correspondencias que llamamos 
\emph{funciones} o \emph{aplicaciones}. 
	
	Cojamos los anteriores conjuntos, $\sA$ y $\sB$. Vamos a decir que el 
conjunto $\sA$ será nuestro conjunto de origen en la correspondencia que será 
una aplicación y a la que llamaremos $\monrel{f}$, y que $\sB$ será nuestro 
conjunto final. Esto lo denotaremos así: $\funct{f} : \sA \fromto \sB$.
	
	
	¿En qué consiste una aplicación? Por decirlo sencillamente, en una serie 
de flechas que van de elementos de $\sA$ a elementos de $\sB$ (una 
correspondencia).
	
	
	Como nuestra intención es comparar conjuntos, impondremos un par de 
condiciones a las correspondencias que nos ayudarán a hacer bien esas 
comparaciones. Las condiciones serán: primero, de cada elemento del conjunto 
origen $\sA$ ha de salir una flecha que ha de dar en el conjunto $\sB$ final. 
Segundo, de cada elemento del conjunto origen, $\sA$ solo puede salir una 
sola flecha al conjunto final $\sB$. Cambiando un poco el lenguaje, cada 
elemento de $\sA$ ha de capturar un elemento de $\sB$, y solo uno. Si todo 
esto ocurre, $\funct{f}$ será una \emph{aplicación} de $\sA$ en $\sB$.
	
	
	
	Ahora bien, la verdad es que hay muchas maneras de establecer la 
aplicación. Algunas más interesantes que otras.
	
	
	La primera forma de establecer la aplicación es fácil de explicar: cada 
par de elementos distintos del conjunto origen, en símbolos, $a b \in \sA 
\;,\; a \neq b$ capturará elementos diferentes del conjunto final, en 
símbolos, $\funct{f}(a) \neq \funct{f}(b)$ (en palabras, si $a$ captura 
mediante $\funct{f}$ y si $b$ captura mediante $\funct{f}$ y si $a \neq b$ 
entonces los elementos capturados son diferentes: los cazadores del conjunto 
inicial no pueden compartir su caza). Pongamos un ejemplo.
	
	
	Vamos a construir una aplicación concreta de $\sA$ a $\sB$. Primero 
comenzaremos por los elementos mencionados $a b \in \sA$.
	
	$$
	\begin{matrix}
		\funct{f}\; :\;   & A & \fromto & B \\
		{}  & a & \mapsto & 2 \\
		{}  & b & \mapsto & 4
		% {}  & c & \mapsto & 7 \\
		% {}  & d & \mapsto & 1 \\
		% {}  & e & \mapsto & 0 \\
	\end{matrix}
	$$
	
	
	Todavía no tenemos una aplicación. Faltan elementos de $\sA$ por capturar 
algún elemento de $\sB$. Pero para $a \in \sA \mapsto \functat{f}{a} = 2 \in 
\sB$, y $b\in \sA \mapsto \functat{f}{b} = 4 \in \sB$, y como $2 \neq 4$ 
tenemos que $a$ y $b$ han capturado elementos distintos ($a \neq b$). Si 
conseguimos seguir en esta tónica, tendremos una aplicación que además tendrá 
la propiedad que buscamos.
	
	
	$$
	\begin{matrix}
		\funct{f}\; :\;   & \sA & \fromto & \sB \\
		{}  & a & \mapsto & 2 \\
		{}  & b & \mapsto & 4 \\
		{}  & c & \mapsto & 7 
		% {}  & d & \mapsto & 1 \\
		% {}  & e & \mapsto & 0 \\
	\end{matrix}
	$$
	
	
	Ahora $\funct{f}(a) \neq \funct{f}(b)$, $\funct{f}(a)=2 \neq 
\functat{f}{c}=7$ y  $\functat{f}{b}=4 \neq \funct{f}{c}=7$ con lo que 
seguimos manteniendo la propiedad que elementos distintos del conjunto 
inicial capturan elementos distintos del conjunto final. Ahora pondré los dos 
elementos que quedan en la colección inicial y conseguiremos que $\funct{f} : 
\sA \fromto \sB$ sea aplicación. Y además con la propiedad buscada.
	
	
	
	$$
	\begin{matrix}
		\funct{f}\; :\;   & \sA & \fromto & \sB \\
		{}  & a & \mapsto & 2 \\
		{}  & b & \mapsto & 4 \\
		{}  & c & \mapsto & 7 \\
		{}  & d & \mapsto & 1 \\
		{}  & e & \mapsto & 0
	\end{matrix}
	$$
	
	
	A la propiedad de las aplicaciones que hemos buscado y que se cumple para 
$\funct{f}$ la llamamos \emph{inyectividad}. Ahora $\funct{f}$ es 
\emph{aplicación inyectiva}.
	
	¿Para qué es interesante una aplicación inyectiva? Fijaos: ahora podemos 
decir que $\sA \preceq \sB$ ($\sA$ es menor o igual que $\sB$). Piénsalo con 
tranquilidad, no hay prisa. Intenta no meter números en tus pensamientos, 
sólo los conjuntos, sus elementos y la aplicación que hemos establecido.
	
	
	Aunque solo sea para ir fijando poco a poco el lenguaje 
lógico-matemático, expondré lo visto en fórmulas:
	
	
	Notación para conjuntos:
	
	
	\begin{flalign*}
		\sA & \;\text{es \textsc{conjunto}} \\
		\sA & = \;\{ a , b , c , d , e \}\\
		& a \in \sA \;\; b \in \sA \;\; c \in \sA \;\; d \in \sA \;\; e \in 
\sA \\
		\sB & \;\text{es \textsc{conjunto}} \\
		\sB & = \;\{ 2 , 4 , 7 , 1 , 0 \}\\
		& 2 \in \sB \;\; 4 \in \sB \;\; 7 \in \sB \;\; 1 \in \sB \;\; 0 \in 
\sB
	\end{flalign*}
	
	
	
	Notación para aplicaciones y aplicaciones inyectivas entre conjuntos:
	
	
	\begin{flalign*}
		\funct{f} \; : \; \sA \fromto \sB \;\text{es \textsc{aplicación}}\; & 
\forall x \in \sA\;\; \exists ! z \in \sB \;\; \functat{f}{x}=z \\
		\funct{f} \; : \; \sA \fromto \sB \;\text{es \textsc{inyectiva}}\; & 
\forall x y \in \sA \; x \neq y \;\Rightarrow\; 
\functat{f}{x}\neq\functat{f}{y}
	\end{flalign*}
	
	
	Tiene que quedar claro que si $\funct{f} : \sA \fromto \sB$ es 
\emph{inyectiva} entonces podemos considerar que $\sA$ puede quedar embutido 
en $\sB$, o dicho de otro modo, en $\sB$ hay suficientes elementos para 
albergar todo $\sA$. Decimos que el \emph{cardinal} de $\sA$ es menor o igual 
que el \emph{cardinal} de $\sB$. El cardinal se denota de varias formas, 
nosotros usaremos el signo numeral $\sharp$ de forma que $\sharp \sA$ es el 
cardinal de $\sA$ y $\sharp \sB$ el cardinal de $\sB$.
	
	
	No importa nada si ahora no te dicen nada las fórmulas, son solo para que 
las vayas viendo.
	
	
	Quedan un par de ideas sobre aplicaciones por explicar. La primera a 
explicar es la de \emph{correspondencia inversa}. Partamos de la que ya 
tenemos hecha:
	
	
	Si $\funct{f}$ es definido:
	$$
	\begin{matrix}
		\funct{f}\; :\;   & \sA & \fromto & \sB \\
		{}  & a & \mapsto & 2 \\
		{}  & b & \mapsto & 4 \\
		{}  & c & \mapsto & 7 \\
		{}  & d & \mapsto & 1 \\
		{}  & e & \mapsto & 0
	\end{matrix}
	$$
	
	
	Entonces definimos la correspondencia inversa de $\funct{f}$ (denotada 
como $\monrelinv{f}$):
	
	$$
	\begin{matrix}
		{\monrelinv{f}}\; :\;   & \sB & \fromto & \sA \\
		{}  & 2 & \mapsto & a \\
		{}  & 4 & \mapsto & b \\
		{}  & 7 & \mapsto & c \\
		{}  & 1 & \mapsto & d \\
		{}  & 0 & \mapsto & e
	\end{matrix}
	$$
	
	
	Es sencillo, si $\funct{f} : \sA \fromto \sB$ es una \emph{aplicación}, 
la inversa es la \emph{correspondencia} ${\monrelinv{f}} : \sB \fromto \sA$ 
dónde se han dado vuelta todas las flechas. Si utilizo la palabra 
\emph{correspondencia} en vez de la palabra función es por un motivo 
sencillo. Para que lo veas voy a definir un nueva aplicación $\funct{g}$ de 
$\sA$ en $\sB$:
	
	
	$$
	\begin{matrix}
		\funct{g}\; :\;   & \sA & \fromto & \sB \\
		{}  & a & \mapsto & 2 \\
		{}  & b & \mapsto & 2 \\
		{}  & c & \mapsto & 7 \\
		{}  & d & \mapsto & 7 \\
		{}  & e & \mapsto & 0 \\
	\end{matrix}
	$$
	
	
	Como puedes ver $\funct{g}$ es una \emph{aplicación}. Cada elemento del 
conjunto origen ha capturado un elemento del conjunto final. Y de cada 
elemento del conjunto inicial sale solo una flecha hacia el final.
	
	
	Ahora vamos a ver como es la \emph{correspondencia inversa} 
${\monrelinv{g}}$:
	
	
	$$
	\begin{matrix}
		{\monrelinv{g}}\; :\;   & \sB & \fromto & \sA \\
		{}  & 2 & \mapsto & a \\
		{}  & 2 & \mapsto & b \\
		{}  & 7 & \mapsto & c \\
		{}  & 7 & \mapsto & d \\
		{}  & 0 & \mapsto & e \\
	\end{matrix}
	$$
	
	
	Puedes observar que ${\monrelinv{g}}$ \emph{no es aplicación}. Hay dos 
elementos del conjunto inicial $\sB$, concretamente $1 4 \in \sB$ que no 
capturan elementos del conjunto final. Pero, además, el elemento $2 \in \sB$ 
captura dos elementos del final $a b \in \sA$, y el elemento $7 \in \sB$ 
captura también dos elementos del final $c d \in \sA$. Las dos condiciones 
que se piden para ser aplicación se incumplen. \emph{Por eso, en general, la 
inversa de una aplicación es una correspondencia}. Pero a veces si que lo es 
como en el caso de ${\funct{f}}$ y ${\monrelinv{f}}$ que son ambas 
aplicaciones. Sin embargo ${\funct{g}}$ es aplicación y ${\monrelinv{g}}$ no 
lo es.
	
	Cuando la inversa de una aplicación es también aplicación, como en el 
caso de ${\funct{f}}$ denotaremos a la inversa como  ${\functinv{f}}$, solo 
para diferenciar esos casos. 
	
	
	El último concepto sobre aplicaciones que vamos a ver, se llama 
\emph{sobreyectividad}. \textbf{\emph{Una aplicación es sobreyectiva cuando 
todos los elementos del conjunto final son capturados}}.
	
	Por ejemplo $\funct{f}$ y ${\functinv{f}}$ son inyectivas y 
sobreyectivas. Sin embargo $\funct{g}$ no es ni inyectiva ni sobreyectiva. En 
este caso ${\monrelinv{g}}$ ni siquiera es aplicación.
	
	
	Retomemos el hilo de cuando un conjunto es más grande que otro: 
\textbf{\emph{si una aplicación cualquiera, digamos}} ${\funct{h}} : {\sX} 
\fromto {\sY}$\textbf{\emph{, es inyectiva decimos que }}${\sX}$ 
\textbf{\emph{es menor o igual (está embutido en) a }}${\sY}$.
	
	Ahora tenemos una nueva forma de ver estas cosas:  \textsf{\emph{si una 
aplicación cualquiera, digamos ${\funct{k}} : {\sV} \fromto {\sW}$, es 
sobreyectiva decimos que $\sV$ es mayor o igual (embute) a $\sW$}}.
	
	Con estas dos sencillas reglas podemos ver cuando dos conjuntos son igual 
de grandes, si ${\funct{h}} : {\sX} \fromto {\sY}$ es a la vez inyectiva y 
sobreyectiva, entonces decimos que $\sX$ y $\sY$ son igual de grandes o 
equipotentes, en símbolos, ${\sX}{\simeq}{\sY}$ o utilizando el símbolo de 
cardinal, $\sharp\sX = \sharp\sY$. Date cuenta que \textbf{\emph{si 
$\funct{h}$ es aplicación inyectiva y sobreyectiva, entonces y solo entonces, 
${\functinv{h}}$ es también aplicación}} (y también inyectiva y 
sobreyectiva).
	
	
	Ya tenemos los criterios para comparar dos colecciones de objetos: dada 
una aplicación entre esas dos colecciones, si esa aplicación es inyectiva, el 
conjunto inicial es menor o igual al final, si es sobreyectiva, el conjunto 
inicial es mayor o igual al final. Finalmente, si esa aplicación es inyectiva 
y sobreyectiva, entonces los dos conjuntos son iguales (equipotentes).




	\emph{\textbf{Aviso a navegantes:} si esa aplicación no es inyectiva ni 
sobreyectiva, esa aplicación no nos dice nada sobre los tamaños de los 
conjuntos.}


	\section{A comparar ...}


	Ha llegado la hora de la verdad. 
	
	
	Vamos a comparar conjuntos mediante funciones. 
	
	
	Eso nos llevará a comprender qué es un número y mucho más...
	
	
	Que entre dos conjuntos $\sA$ y $\sB$ haya un función $\funct{f}:\sA 
\fromto \sB$ dice en principio poco sobre la relación general entre esos 
conjuntos.
	
	
	Pero ¿Qué pasa si decimos que $\funct{f}$ es inyectiva?
	
	
	Recordemos qué quiere decir \emph{función inyectiva}. Decimos que 
$\funct{f} : \sA \fromto \sB$ es \textsc{inyectiva} si para cada dos 
elementos de $a b \in \sA$ tales que $a \neq b$ ocurre que lo que atrapa 
$\funct{f}$ bajo $a$ es distinto de lo que atrapa bajo $b$, $\functat{f}{a} 
\neq \functat{f}{b}$ dónde $\functat{f}{a} \in \sB$ y $\functat{f}{b} \in 
\sB$. Hay suficiente caza en $\sB$ para que cada cazador de $\sA$ encuentre 
su propia presa sin encontrarse con otro cazador. Decimos que $\sA \leqpot 
\sB$ o bien que $\card{\sA} \leq \card{\sB}$. Así que una función inyectiva 
de $\sA$ en $\sB$ embute el conjunto $\sA$ en el conjunto $\sB$ (como se mete 
la carne condimentada en una tripa de cerdo para hacer chorizo). $\sA$ sería 
la carne condimentada y $\sB$ la tripa de cerdo: $\sA$ cabe dentro de $\sB$.


	¿Qué pasa si ahora encontramos una función $\funct{g}:\sB\fromto\sA$ de 
forma que $\funct{g}$ sea \emph{una función inyectiva}?


	Por haber encontrado $\funct{f}$ \emph{inyectiva} ahora sabemos que 
$\sA\leqpot\sB$. Pero como hemos encontrado $\funct{g}$ también 
\emph{inyectiva} y va en el sentido contrario (de $\sB$ a $\sA$), también 
sabemos que $\sB \leqpot \sA$. Si nos dejamos llevar por nuestra intuición 
(en este caso es peligroso) llegamos a que $\card{A}=\card{B}$.


	Como he dicho, es peligroso fiarse de la intuición en estos casos, en que 
tratamos con colecciones (conjuntos), no debidamente definidos. Ya veremos 
que problemas trae esto.
	
	
	Así que... ¿Qué pasa si ahora encontramos una función 
$\funct{h}:\sA\fromto\sB$ de forma que $\funct{h}$ sea \textbf{\emph{una 
función sobreyectiva}}?
	
	
	Primero, antes de comenzar, ¿qué quiere decir que 
$\funct{h}:\sA\fromto\sB$ sea \textbf{\emph{sobreyectiva}}? Recordemos que 
por ser $\funct{h}$ función, todos los elementos del conjunto origen o 
inicial tienen o capturan un elemento del final. La sobreyectividad es una 
propiedad análoga a esta: todos los elementos del conjunto final son 
capturados por algún elemento (y sólo ese elemento) del conjunto inicial. 
Ahora es el conjunto final el que parece capturado dentro del conjunto 
inicial. Cada elemento del conjunto inicial tiene uno del final y no queda 
elemento del final sin capturar. El conjunto final está capturado dentro del 
inicial, cabe dentro del inicial. El inicial es así más grande que el final: 
$\sA \geqpot \sB$ o dicho de otra forma, $\card \sA \geq \card \sB$ o 
también, de forma equivalente, $\numeral \sA \geq \numeral \sB$.



	Por otro lado vamos a tratar la correspondencia inversa de un función 
inyectiva $\funct{f}:\sA \fromto \sB$, que en general será $\monrelinv{f}:\sB 
\fromto \sA$ y será una correspondencia, no necesariamente una función. ¿Cómo 
podemos definir una función que sea inversa de $\funct{f}$? No queda más 
remedio que quitarnos del medio los elementos de $\sB$ que no reciben 
flechas: ${\sB}_{\mthds 1}\subseteq{\sB}$, tal que la función reducida 
${\funct{f}}_{\prime} : \sA \fromto {\sB}_{\mthds 1}$ tal que $ 
{\funct{f}}_{\mthds 1} \left( {x} \right) \defin \funct{f} \left( {x} 
\right)$ y esto $\forall x \in \sA$. $ {\sB}_{\prime} $ se llama la imagen de 
la función $\funct{f}$ denotado en símbolos como $\mathcal{Im}\;\funct{f}$ o 
también $\functat{f}{\sA}$. Pensemos ahora. 
${\funct{f}}_{\prime}:\mathcal{Im}\;\funct{f}\fromto\sA$ dónde 
${\funct{f}}_{\prime}\left(x\right)=\functat{f}{x}$, nos describe una función 
que es a la vez inyectiva y sobreyectiva, \textbf{\emph{una función 
biyectiva}}. Estas son las funciones que al darles la vuelta (encontrar su 
inversa) son también funciones. Pero por ahora eso no importa mucho. Lo que 
sí que sabemos es que ${\funct{f}}_{\prime}^{-1} : \mathcal{Im} \; \funct{f} 
\fromto \sA$ es función. Hemos hecho la trampa de quitar de enmedio los 
elementos de $\sB$ huérfanos de elemento inicial. Ahora podemos ampliar 
${\funct{f}}_{\prime}^{-1}$ de forma que vaya de $\sB$ a $\sA$ como función. 
Es fácil: supongamos que cualquier elemento de $\sB$ que no esté en la imagen 
por $\funct{f}$ de $\sA$, esto $\forall y \in \sB \setminus \mathcal{Im} \; 
\funct{f}$ lo aplicamos a un elemento elegido cualquiera de $sA$, por ejemplo 
$\exists x_0 \in \sA$, y ahora aplicamos $\forall y \in \sB \setminus 
\mathcal{Im} \; \funct{f} \quad {{\funct{f}}_{\prime\prime}}^{ -1 } \left( y 
\right) \defin x_0$. Todo lo demás queda igual. Ahora 
${{\funct{f}}_{\prime\prime}}^{ -1 } \circ \funct{f} : \sA \fromto \sA = { 
\funct{1} }_{ \sA }$.



	Creo que he ido un poco rápido en el último párrafo. Intentaré 
solucionarlo en éste. Supongamos que tenemos un par de conjuntos, primero es 
el conjunto vacío, $ \emptyset $, y el segundo será un conjunto con un solo 
elemento, $ \sA \defin \set{*} $. Intentemos ver cuáles son las posibles 
funciones de $ \emptyset \fromto \set{*} $. Pero esto es sencillo: como en el 
conjunto inicial no hay elementos, no puede salir ninguna flecha. Así que el 
conjunto de funciones de $ \emptyset \fromto \set{*} $ no está vacío, tiene 
un solo elemento $\funct{f} : \emptyset \fromto \set{*} :: \textsc{sin 
asignaciones}$. Date cuenta que alrevés, $\set{*} \fromto \emptyset$ no 
funciona, el conjunto de estas funciones es vacío, ya que si no pueden salir 
flechas del conjunto inicial, y en el conjunto inicial hay al menos un 
elemento, el que no haya flechas quiere decir que son correspondencias sin 
flechas, pero no funciones. Así que la función antes hallada $\funct{f}$ no 
tiene inversa, y no es biyectiva. Esto es $\monrelinv{f}$ no es función. 
$\funct{f}$ es inyectiva pero no es sobreyectiva. Así sabemos que $\emptyset 
\precnsim \set{*}$ y en este caso incluso $\emptyset \subsetneq \set{*}$. Si 
\textbf{\emph{nombramos}} el cardinal del vacío por $0 \defin 
\card{\emptyset}$ y si también \textbf{\emph{nombramos}} $1 \defin 
\card{\set{*}}$ tenemos que $0 \lneqq 1$. Ahora tienes que darte cuenta que 
dónde hemos puesto $\set{*}$ podíamos haber puesto cualquier otro conjunto 
con uno o más elementos. Todos los razonamientos sería idénticos. Por lo 
tanto, $0$ es el más pequeño de los cardinales posibles. De otro lado todo 
conjunto es equipotente consigo mismo: $\emptyset \simeq \emptyset$. También 
es aquí conveniente que te des cuenta de otra verdad: solo existe un conjunto 
vacío. Dos conjuntos no se distinguen más que por sus elementos, y como el 
conjunto vacío no tiene elementos, no puede variar de ninguna forma: hay un 
solo conjunto vacío. Pongamos que hay dos conjuntos vacíos 
$\emptyset^\prime\text{ y }\emptyset^{\prime\prime}$. Como $\forall x \in 
\emptyset^\prime\quad x \in \emptyset^{\prime\prime}$, y viceversa, $\forall 
y \in \emptyset^{\prime\prime}\quad y \in \emptyset^\prime$ no queda más que 
$\emptyset^\prime = \emptyset^{\prime\prime} = \emptyset$.


	Seguimos con el párrafo anterior. Ahora vamos a ver los conjuntos de un 
solo elemento con conjuntos de, bien de uno, bien de dos elementos. El 
conjunto inicial será $\set{*}$ y el final $\set{\clubsuit,\diamondsuit}$. 
Veamos cuántas funciones podemos hacer:


\begin{flalign*}
  & \funct{f}_1 \quad : \quad \set{*} \fromto \quad 
  \set{\clubsuit,\diamondsuit} &\\
  & {\quad\quad} \quad :: \quad\quad * \mapsto \quad \clubsuit &\\
  & {}           &\\
  & \funct{f}_2 \quad : \quad \set{*} \fromto \quad 
  \set{\clubsuit,\diamondsuit} &\\
  & {\quad\quad} \quad :: \quad\quad * \mapsto \quad \diamondsuit &\\
\end{flalign*}


	Ahora veamos las relaciones reversas de éstas:


\begin{flalign*}
  & \monrelinv{f}_1 \quad : \quad \set{\clubsuit,\diamondsuit} \fromto \quad 
  \set{*} &\\
  & {\quad\quad} \quad :: \quad\quad \clubsuit \mapsto \quad * &\\
  & {\quad\quad} \quad :: \quad\quad \diamondsuit \quad \quad\quad          
  &\\
  & \monrelinv{f}_2 \quad : \quad \set{\clubsuit,\diamondsuit} \fromto \quad 
  \set{*} &\\
  & {\quad\quad} \quad :: \quad\quad \diamondsuit \mapsto \quad * &\\
  & {\quad\quad} \quad :: \quad\quad \clubsuit \quad \quad\quad          &\\
\end{flalign*}


	Es claro que ni $\monrelinv{f}_1$ ni $\monrelinv{f}_2$ son funciones (no 
cumplen que cada elemento inicial capture uno del final). Por tanto podemos 
ver que tanto $\funct{f}_1$ como $\funct{f}_2$ son funciones inyectivas y son 
no sobreyectivas. No son funciones biyectivas. Llegamos a la conclusión que 
$\set{*} \precnsim \set{\clubsuit,\diamondsuit}$. Si ahora 
\textbf{\emph{llamamos}} a $2 \defin \card{\set{\clubsuit,\diamondsuit}}$ es 
claro que $1 \lneq 2$.


	A partir de aquí es fácil ver algunas propiedades curiosas: por ejemplo, 
las funciones entre $\set{*}$ y $\set{\diamondsuit}$ son solo una, no puede 
haber otra, y que además esa única función es necesariamente biyectiva: 
$\functat{f}{*}={\diamondsuit}$ y $\functinvat{f}{\diamondsuit}={*}$.


	Así que diremos (muy informalmente como ya veremos) que las funciones de 
un conjunto de elemento en otro de un solo elemento son siempre biyectivas y 
por lo tanto todos los conjuntos de un solo elemento tienen el mismo 
cardinal, a saber, el que hemos definido como $1$. Además $1$ es 
estrictamente menor que cualquier otro cardinal que no sea $0$ (cardinal de 
$\emptyset$) o $1$.


	Sigamos con los conjuntos tipo $\set{\diamondsuit,\clubsuit}$. Supongamos 
que tenemos otro conjunto $\set{\heartsuit,\spadesuit}$. Veamos las funciones 
de $\set{\diamondsuit,\clubsuit}\fromto\set{\heartsuit,\spadesuit}$.


\begin{flalign*}
& \funct{f}_1 : \set{\diamondsuit, \clubsuit} \fromto 
\set{\heartsuit,\spadesuit} &\\
& \quad :: \quad \diamondsuit \quad \mapsto \quad \heartsuit &\\
& \quad :: \quad \clubsuit \quad \mapsto \quad \spadesuit &\\
& &\\
& \funct{f}_2 : \set{\diamondsuit, \clubsuit} \fromto 
\set{\heartsuit,\spadesuit} &\\
& \quad :: \quad \diamondsuit \quad \mapsto \quad \spadesuit &\\
& \quad :: \quad \clubsuit \quad \mapsto \quad \heartsuit &\\
& &\\
& \funct{f}_3 : \set{\diamondsuit, \clubsuit} \fromto 
\set{\heartsuit,\spadesuit} &\\
& \quad :: \quad \diamondsuit \quad \mapsto \quad \heartsuit &\\
& \quad :: \quad \clubsuit \quad \mapsto \quad \heartsuit &\\
& &\\
& \funct{f}_4 : \set{\diamondsuit, \clubsuit} \fromto 
\set{\heartsuit,\spadesuit} &\\
& \quad :: \quad \diamondsuit \quad \mapsto \quad \spadesuit &\\
& \quad :: \quad \clubsuit \quad \mapsto \quad \spadesuit &
\end{flalign*}


	Y sus inversas:


\begin{flalign*}
& \functinv{f}_1 : \set{\heartsuit, \spadesuit} \fromto 
\set{\diamondsuit,\clubsuit} &\\
& \quad :: \quad \heartsuit \quad \mapsto \quad \diamondsuit &\\
& \quad :: \quad \spadesuit \quad \mapsto \quad \clubsuit &\\
& &\\
& \functinv{f}_2 : \set{\heartsuit, \spadesuit} \fromto 
\set{\diamondsuit,\clubsuit} &\\
& \quad :: \quad \spadesuit \quad \mapsto \quad \diamondsuit &\\
& \quad :: \quad \heartsuit \quad \mapsto \quad \clubsuit &\\
& &\\
& \monrelinv{f}_3 : \set{\heartsuit, \spadesuit} \fromto 
\set{\diamondsuit,\clubsuit} &\\
& \quad :: \quad \heartsuit \quad \mapsto \quad \diamondsuit &\\
& \quad :: \quad \heartsuit \quad \mapsto \quad \clubsuit &\\
& \quad :: \quad \spadesuit \quad \quad \quad \quad &\\
& &\\
& \monrelinv{f}_4 : \set{\heartsuit, \spadesuit} \fromto 
\set{\diamondsuit,\clubsuit} &\\
& \quad :: \quad \spadesuit \quad \mapsto \quad \diamondsuit &\\
& \quad :: \quad \spadesuit \quad \mapsto \quad \clubsuit &\\
& \quad :: \quad \heartsuit \quad \quad \quad \quad &
\end{flalign*}


	Podemos ver que $\funct{f}_1$ y $funct{f}_2$ son funciones biyectivas, 
esto es, a cada elemento del conjunto inicial le corresponde uno del final, 
distinto para cada uno (\emph{inyectividad}) y no queda ningún elemento del 
conjunto final sin quedar capturado (ninguno queda sin elemento origen, 
\emph{sobreyectividad}). Por ello, tanto $\functinv{f}_1$ como  
$\functinv{f}_2$ son también funciones y además biyectivas.


	Lo que realmente cuenta es la existencia de alguna función biyectiva. Aún 
así podemos ver que $\funct{f}_3$ y $\funct{f}_4$ son funciones, pero no nos 
interesan desde que ninguna de ellas es ni inyectiva ni sobreyectiva.


	Por la existencia de las biyectividades $\funct{f}_1$ o $\funct{f}_2$ 
podemos decir que estos dos conjuntos tienen la misma cardinalidad, 
$\set{\heartsuit, \spadesuit} \simeq \set{\diamondsuit,\clubsuit}$, o de otra 
forma $\card{\set{\heartsuit, \spadesuit}} = 
\card{\set{\diamondsuit,\clubsuit}} = 2$. Podemos llamar $2$ a todas las 
colecciones que sean equipotentes a las anteriores.


	Como puedes ver, este proceso puede seguir, haciendo crecer los conjuntos 
con nuevos elementos y obtendríamos nuevos cardinales. Parece que podemos ver 
cada cardinal como la colección de colecciones equipotentes entre sí.


	¿Hasta dónde puede llevarnos este proceso?


	Podríamos seguir con un conjunto $\set{\nabla , \triangle, \square}$ y 
otro $\set{\clubsuit,\diamondsuit,\spadesuit}$. Las únicas biyecciones entre 
estos dos conjuntos son en total seis (se llaman permutaciones):
\begin{flalign*}
  & \funct{f}_1 : \set{\nabla , \triangle, \square} \fromto 
  \set{\clubsuit,\diamondsuit,\spadesuit}&\\
  & {\quad} :: {\quad\quad} \nabla {\quad\quad} \mapsto {\quad\quad} 
  \clubsuit &\\
  & {\quad} :: {\quad\quad} \triangle {\quad\quad} \mapsto {\quad\quad} 
  \diamondsuit &\\
  & {\quad} :: {\quad\quad} \square {\quad\quad} \mapsto {\quad\quad} 
  \spadesuit &
\end{flalign*}
\begin{flalign*}
  & \funct{f}_2 : \set{\nabla , \triangle, \square} \fromto 
  \set{\clubsuit,\diamondsuit,\spadesuit}&\\
  & {\quad} :: {\quad\quad} \nabla {\quad\quad} \mapsto {\quad\quad} 
  \clubsuit &\\
  & {\quad} :: {\quad\quad} \triangle {\quad\quad} \mapsto {\quad\quad} 
  \spadesuit &\\
  & {\quad} :: {\quad\quad} \square {\quad\quad} \mapsto {\quad\quad} 
  \diamondsuit &
\end{flalign*}
\begin{flalign*}
  & \funct{f}_3 : \set{\nabla , \triangle, \square} \fromto 
  \set{\clubsuit,\diamondsuit,\spadesuit}&\\
  & {\quad} :: {\quad\quad} \nabla {\quad\quad} \mapsto {\quad\quad} 
  \diamondsuit &\\
  & {\quad} :: {\quad\quad} \triangle {\quad\quad} \mapsto {\quad\quad} 
  \clubsuit &\\
  & {\quad} :: {\quad\quad} \square {\quad\quad} \mapsto {\quad\quad} 
  \spadesuit &
\end{flalign*}
\begin{flalign*}
  & \funct{f}_4 : \set{\nabla , \triangle, \square} \fromto 
  \set{\clubsuit,\diamondsuit,\spadesuit}&\\
  & {\quad} :: {\quad\quad} \nabla {\quad\quad} \mapsto {\quad\quad} 
  \diamondsuit &\\
  & {\quad} :: {\quad\quad} \triangle {\quad\quad} \mapsto {\quad\quad} 
  \spadesuit &\\
  & {\quad} :: {\quad\quad} \square {\quad\quad} \mapsto {\quad\quad} 
  \clubsuit &
\end{flalign*}
\begin{flalign*}
  & \funct{f}_5 : \set{\nabla , \triangle, \square} \fromto 
  \set{\clubsuit,\diamondsuit,\spadesuit}&\\
  & {\quad} :: {\quad\quad} \nabla {\quad\quad} \mapsto {\quad\quad} 
  \spadesuit &\\
  & {\quad} :: {\quad\quad} \triangle {\quad\quad} \mapsto {\quad\quad} 
  \diamondsuit &\\
  & {\quad} :: {\quad\quad} \square {\quad\quad} \mapsto {\quad\quad} 
  \clubsuit &
\end{flalign*}
\begin{flalign*}
  & \funct{f}_6 : \set{\nabla , \triangle, \square} \fromto 
  \set{\clubsuit,\diamondsuit,\spadesuit}&\\
  & {\quad} :: {\quad\quad} \nabla {\quad\quad} \mapsto {\quad\quad} 
  \spadesuit &\\
  & {\quad} :: {\quad\quad} \triangle {\quad\quad} \mapsto {\quad\quad} 
  \clubsuit &\\
  & {\quad} :: {\quad\quad} \square {\quad\quad} \mapsto {\quad\quad} 
  \diamondsuit &
\end{flalign*}


	Como lo único que nos interesa es que haya una biyección, no volveremos a 
hacer estas largas listas (conforme crezca el conjunto en elementos más 
crecerá la lista de biyecciones, y en general de funciones).


	Sin explicitarlo totalmente supongamos que ahora nuestros conjuntos son, 
como conjunto original $\set{\nabla , \triangle, \square}$ y como conjunto 
final $\set{\clubsuit,\diamondsuit,\spadesuit,\heartsuit}$. Podemos ver que 
hay funciones inyectivas del conjunto origen en el final. Pero podemos ver 
también que no hay ninguna función sobreyectiva del original en el final. 
Esto es, podemos ver que $\set{\nabla , \triangle, \square} \precnsim 
\set{\clubsuit,\diamondsuit,\spadesuit,\heartsuit}$. Si ahora cambiamos los 
papeles, esto es, $\set{\clubsuit,\diamondsuit,\spadesuit,\heartsuit}$ será 
ahora el conjunto inicial y $\set{\nabla , \triangle, \square}$ será ahora el 
final, podemos ver fácilmente que hay funciones sobreyectivas pero que 
ninguna es inyectiva. Por lo tanto, $\card{\set{\nabla , \triangle, \square}} 
\lneq \card{\set{\clubsuit,\diamondsuit,\spadesuit,\heartsuit}}$, esto es, 
menor y no igual, que se suele decir estrictamente menor. Además es fácil ver 
que si añadimos nuevos elementos al conjunto mayor (el de los palos de la 
baraja) la situación seguirá siendo la misma, el conjunto de figuras 
geométricas seguirá siendo estrictamente menor que el de palos de la baraja 
aumentada. Podríamos llamar a este cardinal $3$. Esto es $3 \defin 
\card{\set{\nabla , \triangle, \square}}$. Lo podemos ver como la colección 
(o clase) de los conjuntos equipotentes (con biyecciones) al actual conjunto 
de formas geométricas.


	Esta claro que este proceso podríamos llevarlo \emph{ad infinitum}. 
Acabamos de reproducir los números normales y corrientes (los naturales, que 
se suelen representar por $\sN$), con la pequeña diferencia que hemos 
introducido el número (cardinal) $0 \defin \card{\emptyset}$, que no se suele 
meter en ese conjunto.


	Podríais decirme que para este viaje no necesitábamos estas alforjas. Y 
tendríais razón si solo nos limitáramos a los conjuntos pequeños (o finitos).


	Pero antes de seguir adelante podemos intentar ver un poco mejor qué es 
lo que hemos visto.


	Dadas dos colecciones, $\sA$ y $\sB$, un número natural sería (con lo 
visto hasta ahora), el cardinal de un conjunto. Por ejemplo $a \defin 
\card{\sA}$ y $b \defin \card{\sB}$. Veamos las relaciones que hemos visto 
hasta ahora.

\begin{thm}
  Dados el conjunto $\sA$ con cardinal $a \defin \card{\sA}$ y el conjunto 
  $\sB$ con cardinal $b \defin \card{\sB}$, la relación de orden o igualdad 
  entre $a$ y $b$ vienen dadas por:
  \begin{enumerate}
  	\item Si existe una biyección de $\sA \fromto \sB$, entonces decimos 
  que $\sA \simeq \sB$ y por lo tanto que $a = b$.
    \item Si no existe una biyección entre $\sA$ y $\sB$, entonces, 
    claramente, $a\neq b$.
        \begin{enumerate}
        \item Si existe una función inyectiva de $\sA \fromto \sB$, 
        entonces, como no existen biyecciones, tampoco puede haber 
        sobreyecciones de $\sA \fromto \sB$. Por la misma causa tenemos 
        que no hay inyecciones de $\sB \fromto \sA$ (esto, como estamos 
        en matemáticas, habría que probarlo). $\sA \precnsim \sB$ y $a = 
        \card{\sA}\lneq \card{\sB} = b$.
        \item Si existe una función sobreyectiva de $\sA \fromto \sB$, 
        entonces, como no existen biyecciones, tampoco puede haber 
        inyecciones de $\sA \fromto \sB$. Por la misma causa tenemos que 
        no hay sobreyecciones de $\sB \fromto \sA$ (de nuevo, habría que 
        probarlo). $\sB \precnsim \sA$ y $b = \card{\sB}\lneq \card{\sA} 
        = a$.
        \item No puede ocurrir que dados dos conjuntos entre los que se 
        puede establecer algunas funciones, todas ellas sean a la vez no 
        inyectivas y no sobreyectivas. Esto es, hay un orden completo 
        entre los conjuntos (esto también hay que demostrarlo). Esto 
        quiere decir que $a \lneq b \vee a \gneq b \vee a = b$.
       \end{enumerate}
  \end{enumerate}
\end{thm}


	Profundicemos algo más en lo que hemos estado haciendo todo este largo 
capítulo de comparaciones.


	Al principio partimos de un conjunto sin elementos $\emptyset \defin 
\set{}$. Después añadimos un elemento. Esto se hace con una operación que se 
llama \emph{unión}:

$$
\set{\diamondsuit} = \set{} \cup \set{\diamondsuit}
$$


	Por ahora esta operación entre conjuntos no parece muy útil. Pero lo 
importante es que juntamos todos los elementos de los dos conjuntos en uno 
solo. Así $\emptyset \cup \emptyset = \emptyset$. Es claro porque no hay 
elementos que unir. De hecho, dado un conjunto cualquiera $\sA$, la unión con 
el vacío no hace nada, $\emptyset \cup \sA = \sA$. Es como sumar cero, $1+0 = 
1$ y $0 + 163749 = 163749$. En este caso concreto, cuando uno de los dos 
conjuntos es vacío, el cardinal es la suma aritmética de los cardinales, y 
coincide con el cardinal de la unión de los conjuntos.


	Avancemos a otro caso. Tenemos el conjunto $\sA \defin 
\set{\diamondsuit}$ y $\sB \defin \set{\clubsuit}$. Si queremos hacer la 
unión de los dos conjuntos nos quedará:


$$
\sA\cup\sB = \set{\diamondsuit}\cup\set{\clubsuit} = 
\set{\diamondsuit,\clubsuit}
$$


	En este caso es claro que $1 \defin \card{\sA}$ y $1 \defin \card{\sB}$, 
y $2 \defin \card{\sA\cup\sB} = \card{\sA}+\card{\sB} = 1 + 1$.


	Podríamos pensar que $\card{\sA\cup\sB}=\card{\sA}+\card{\sB}$: 
¡\textbf{\emph{Pero es falso}}! Pongamos un contraejemplo simple: $\card{ 
\set{ \triangle, \square} \cup \set{ \triangle, \bigcirc}} = \card{ \set{ 
\triangle,\square,\bigcirc} }=3\neq 2 + 2$.


	La verdad es que el contraejemplo es aún más fácil: $\card{ \set{*} \cup 
\set{*}} = \card{ \set{*} } = 1 \neq 1 + 1$.


	El problema son los elementos que se repiten. Para solucionar este 
problema tendremos que ver otra operación entre conjuntos: la que nos da los 
elementos repetidos. A esta operación la llamamos intersección.


	La intersección de dos conjuntos $\sA\cap\sB$ es el conjunto de elementos 
comunes o repetidos de esos conjuntos. Por ejemplo, $\set{*}\cap\set{*} = 
\set{*}$. Otro ejemplo, $\emptyset \cap \sA =  \emptyset$, porque no puede 
repetirse ningún elemento cuando una de las partes es el vacío. Aún más, como 
todo conjunto tiene todos sus elementos repetidos consigo mismo $\sA \cap \sA 
= \sA$.


	¿Cómo nos soluciona esta operación el problema anterior? Es fácil, si 
unimos dos conjuntos sale una suma de cardinales que puede ser mayor que el 
cardinal de la unión. Porque los elementos pueden estar repetidos. Pues lo 
único que tendremos que hacer es quitar el cardinal de los elementos 
repetidos:


$$
\card{\sA \cup \sB} = \card{\sA} + \card{\sB} - \card{\sA \cap \sB}
$$


	O dicho de otra manera:


$$
\card{\sA\cup\sB}=\card{\sA}+\card{\sB}\quad \Leftrightarrow\quad \sA \cap 
\sB = \emptyset
$$


	Otra operación entre conjuntos interesante es la resta de conjuntos: 
$\sA\setminus\sB$ es el conjunto de elementos de $\sA$ que \textsc{no} están 
en $\sB$.


	Una forma técnica de describir lo anterior es:

$$
\sA \setminus \sB \quad \defin \quad \setprop {x\in\sA} {x\notin\sB}
$$


	De camino nos cuesta poco hacer lo mismo con la intersección:


$$
\sA \cap \sB \quad \defin \quad \setprop{x\in\sA}{x\in\sB}
$$



	Y lo mismo para la unión:


$$
\sA \cup \sB \quad \defin \quad \setprop{x}{x\in\sA \quad \text{\textsc {ó} } 
\quad x\in\sB}
$$


\paragraph{}


Veamos en qué punto estamos. Pongamos una lista de lo que hemos explicado:
\begin{description}
  \item[1] Hemos visto qué sea un conjunto: una colección de elementos.
  \item[2] Que un conjunto se diferencia de otro porque tiene elementos 
  distintos. Un conjunto no tiene elementos repetidos porque no podría 
  diferenciarse del conjunto con elementos repetidos.
  \item[3] Hemos visto las relaciones entre conjuntos. Hay que decir cual 
  es el conjunto inicial y cual es el final, y decir como capturan los 
  elementos del conjunto inicial a elementos del conjunto final. No hay 
  ninguna restricción para las relaciones.
  \item[4] Como nuestro cometido es comparar conjuntos en tanto si son más 
  grandes o más pequeños o iguales que otros, hemos introducido unas 
  restricciones a las relaciones que nos hacen conseguir funciones. Las 
  funciones son relaciones que en su conjunto inicial todos los elementos 
  cazan elementos del final y además se les pide que solo cacen uno.
  \item[5] Aún con las restricciones dichas, no todas las funciones nos van 
  a servir para comparar conjuntos. Si la función es tal que dos elementos 
  (distintos) del conjunto inicial capturan elementos distintos del 
  conjunto final, decimos que la función es inyectiva. De otro lado, cuando 
  la función consigue no haya elementos del conjunto final sin capturar, la 
  función es sobreyectiva.
  \item[6] Si una función es inyectiva, entonces el conjunto inicial es 
  menor o igual (de grande) que el conjunto final.
  \item[7] Si una función es sobreyectiva, entonces, el conjunto inicial es 
  mayor o igual (de grande) que el conjunto final.
  \item[8] Si una función es inyectiva y sobreyectiva, la función se llama 
  biyectiva, y nos indica que los dos conjuntos son igual de grandes. Las 
  funciones biyectivas tienen sus relaciones invertidas que también son 
  funciones, la función inversa. Y esta función inversa es de nuevo 
  biyectiva.
  \item[9] Si existe una biyección entre dos conjuntos (inicial y final), 
  decimos que ambos conjuntos tienen el mismo cardinal. Podríamos decir que 
  todos los conjuntos que tienen biyecciones entre si son el cardinal de 
  cualquiera de ellos.
  \item[10] Si no existe biyección entre dos conjuntos (inicial y final), 
  pero existe una función inyectiva entre ellos, el cardinal del inicial es 
  estrictamente menor que el cardinal del final. No puede haber (se cumple) 
  en este caso, funciones sobreyectivas.
  \item[11] Si no existe biyección entre dos conjuntos (inicial y final), 
  pero existe una función sobreyectiva entre ellos, el cardinal del inicial 
  es estrictamente mayor que el cardinal del final. No puede haber (se 
  cumple) en este caso, funciones inyectivas.
  \item[12] Para los conjuntos pequeños (que no tienen una infinidad de 
  elementos, como por ejemplo los números), podemos ver que sus cardinales 
  coinciden con los números naturales incluido el cero.
  \item[13] Para ver el punto anterior hemos desarrollado una operación 
  entre conjuntos que se llama unión. Vemos que para obtener conjuntos más 
  grandes nos podemos servir de la unión, pero que ésta no siempre los hace 
  crecer. A veces los deja igual.
  \item[14] Hemos visto otra operación entre conjuntos: la intersección. 
  Los elementos comunes, que nos da un conjunto o menor o igual que los 
  originales.
  \item[15] Aunque no vale directamente para lo que pretendemos, hemos 
  introducido la operación entre conjuntos de diferencia o resta 
  conjuntista.
\end{description}


\section{ ¿Se pueden comparar conjuntos realmente grandes? }


	Hasta el momento hemos comparado conjuntos pequeños. Pero ¿qué me decís 
de un conjunto como $\setN$ que decimos que es infinito y por ejemplo el 
conjunto de los naturales pares, llamémosle $\textsf{Par}$. Cada número par 
podemos ponerlo como $2 \cdot n \in \textsf{Par}$ siendo $n \in \setN$. Dicho 
de la manera formal que introduje al final de la sección anterior: 
$\textsf{Par} \defin \setprop{2 \cdot n}{n \in \setN}$. Evidentemente, 
$\textsf{Par} \subsetneq \setN$. Está dentro porque los números pares son 
números naturales como los demás. Y no son el mismo conjunto, porque $\setN$ 
tiene muchos elementos, todos los impares $\set{1, 3, 5, 7, 9, ...}$ que no 
están en los naturales.


	Lo primero que se os puede venir a la cabeza es: pero si $\textsf{Par}$ 
está completamente dentro de $\setN$ y además no tiene, ni mucho menos, todos 
los elementos de $\setN$... ¿No debería ser el cardinal de los pares menor 
que el de los naturales? Aquí desde luego no nos vale con contar como 
haríamos habitualmente con conjuntos pequeños. ¿Que número utilizamos como el 
cardinal?


	Bien. Esto va a ser grande. Os digo que hay una biyección entre ambos 
conjuntos. Os lo voy a demostrar:


\begin{flalign*}
  \funct{f} :\; & \textsf{Par} \fromto \setN & \\
  {\phantom{\funct{f} :}}& \; 2\cdot n \; \mapsto \; n &
\end{flalign*}


	Lo primero que afirmo es $\funct{f}$ es una \textbf{\emph{función}}. No 
hay número par que se quede sin su número corriente. Y todo número par atrapa 
exactamente un número corriente y ninguno más. Vamos a deletrear la función 
$\funct{f}$ un poco más:


\begin{flalign*}
  \funct{f} :\; & \textsf{Par} \fromto \setN & \\
  {\phantom{\funct{f} :}}& \; 0 = 2 \cdot 0 \; \mapsto \; 0 &\\
  {\phantom{\funct{f} :}}& \; 2 = 2 \cdot 1 \; \mapsto \; 1 &
\end{flalign*}
\begin{flalign*}
  {\phantom{\funct{f} :}}& \; 4 = 2 \cdot 2 \; \mapsto \; 2 &\\
  {\phantom{\funct{f} :}}& \; 6 = 2 \cdot 3 \; \mapsto \; 3 &\\
  {\phantom{\funct{f} :}}& \; 8 = 2 \cdot 4 \; \mapsto \; 4 &
\end{flalign*}
\begin{flalign*}
  {\phantom{\funct{f} :}}& \;10 = 2 \cdot 5 \; \mapsto \; 5 &\\
  {\phantom{\funct{f} :}}& \;12 = 2 \cdot 6 \; \mapsto \; 6 &\\
  {\phantom{\funct{f} :}}& \;14 = 2 \cdot 7 \; \mapsto \; 7 &
\end{flalign*}
\begin{flalign*}
  {\phantom{\funct{f} :}}& \;16 = 2 \cdot 8 \; \mapsto \; 8 &\\
  {\phantom{\funct{f} :}}& \;18 = 2 \cdot 9 \; \mapsto \; 9 &
\end{flalign*}
\begin{flalign*}
  {\phantom{\funct{f} :}}& ... &\\
  {\phantom{\funct{f} :}}& \quad n_{par} \; \mapsto \; \frac{n_{par}}{2}&
\end{flalign*}


	Lo primero es: ¿no es verdad que todo número par es divisible por $2$ de 
manera exacta (con resto $0$ o, lo que es lo mismo, sin resto)? Entonces para 
cada número par existe un elementos de los naturales que es el par dividido 
de forma exacta por dos. Y como esto ocurre para todos los pares no nos queda 
más remedio que decir que la relación antes descrita \textbf{\emph{es una 
función}}. \textbf{\textsf{Quod erat demostrandum}}.


	La parte más sorprendente es que en la columna de los naturales parece 
que aparecen todos los naturales. La afirmación que se sigue es que la 
función $\funct{f}$ es \textbf{\emph{función sobreyectiva}}. Cogemos un 
número cualquiera de $\setN$, digamos $n$, que está en el conjunto final. 
Tendrá un elemento de los pares que lo atrape? Veamos, si multiplicamos $2 
\cdot n$ nos sale un número par. Siempre podemos multiplicar por dos. Lo que 
afirmo es que $n_{par} \defin 2 \cdot n$ es el número origen que atrapa al 
número $n$ elegido (cualquiera puede ser). Efectivamente, $\frac{n_{par}}{2} 
= \frac{2 \cdot n}{2} = \frac{n}{1} = n$. \textbf{\textsf{Quod erat 
demostrandum}}.

	Pero la afirmación que voy a hacer ahora también es válida. $\funct{f}$ 
es una \textbf{\emph{función inyectiva}}. Si no te convence simplemente la 
fórmula: dos números que se dividen de forma exacta entre dos y que son 
distintos entre sí, al dividirlos por dos, de forma exacta darán dos números 
diferentes como cocientes. Pongamos que no lo terminas de ver. Así que vamos 
a suponer que para dos números distintos $n_{par}\neq m_{par}$ del conjunto 
original obtiene el mismo número natural $p$ como imagen. Entonces al 
multiplicar ese número por dos nos ha de dar los dos pares: $2 \cdot p = 
n_{par}$ y $2 \cdot p = m_{par}$. La conclusión es $n_{par} = 2 \cdot p = 
m_{par}$ y de ahí que $n_{par} = m_{par}$. Pero partíamos de $n_{par} \neq 
m_{par}$. Conclusión primera: $n_{par} = m_{par} \quad \textsc{y} \quad 
n_{par} \neq m_{par}$. Más claro aún: $\textsf{SI}\;\; n_{par} = m_{par} 
\quad \textsc{y} \quad \textsf{NO}\;\; n_{par} = m_{par}$. \textsc{Absurdo}. 
\textsc{Luego no existe }$p\in\textsf{Par}$ tal que 
$p=\functat{f}{n_{par}}=\functat{f}{m_{par}}$. Y la función $\funct{f}$ es 
\textbf{\emph{función inyectiva}}. \textsf{Quod erat demostrandum}.



	Por último, al ser $\funct{f}$ función inyectiva y sobreyectiva 
encontramos que la definición de función biyectiva. $\funct{f}$ es 
\textbf{\emph{función biyectiva}}. \textbf{\textsf{Quod erat demostrandum}}.


	Algunos dirán: pero decíamos que las funciones biyectivas tiene relación 
reversa que es función y es además biyectiva. Y llevan razón. Pongo a 
continuación la función inversa de $\funct{f}$:


\begin{flalign*}
  \functinv{f} :\; & \setN \fromto \mathsf{Par} & \\
  {\phantom{\funct{f} :}}& \; 0 \; \mapsto \; 0 = 2 \cdot 0 &
\end{flalign*}
\begin{flalign*}
  {\phantom{\funct{f} :}}& \; 1 \; \mapsto \; 2 = 2 \cdot 1 &\\
  {\phantom{\funct{f} :}}& \; 2 \; \mapsto \; 4 = 2 \cdot 2 &\\
  {\phantom{\funct{f} :}}& \; 3 \; \mapsto \; 6 = 2 \cdot 3 &
\end{flalign*}
\begin{flalign*}
  {\phantom{\funct{f} :}}& \; 4 \; \mapsto \; 8 = 2 \cdot 4 &\\
  {\phantom{\funct{f} :}}& \; 5 \; \mapsto \;10 = 2 \cdot 5 &\\
  {\phantom{\funct{f} :}}& \; 6 \; \mapsto \;12 = 2 \cdot 6 &
\end{flalign*}
\begin{flalign*}
  {\phantom{\funct{f} :}}& \; 7 \; \mapsto \;12 = 2 \cdot 7 &\\
  {\phantom{\funct{f} :}}& \; 8 \; \mapsto \;12 = 2 \cdot 8 &\\
  {\phantom{\funct{f} :}}& \; 9 \; \mapsto \;12 = 2 \cdot 9 &
\end{flalign*}
\begin{flalign*}
  {\phantom{\funct{f} :}}& ... &\\
  {\phantom{\funct{f} :}}& \quad n \; \mapsto \; n_{par} = 2 \cdot n&
\end{flalign*}


	Vamos a componer $\funct{f}\circ\functinv{f} : \setN \fromto \textsf{Par} 
\fromto \setN$.


\begin{flalign*}
  \funct{f}\circ\functinv{f} :\; & \setN \quad \fromto \quad \mathsf{Par} 
  \quad \fromto \quad \setN& \\
  {\phantom{\funct{f} :}}& \; 0 \; \mapsto \; 0 = 2 \cdot 0  \; \mapsto \; 0 
  = \frac{0}{2} &
\end{flalign*}
\begin{flalign*}
  {\phantom{\funct{f} :}}& \; 1 \; \mapsto \; 2 = 2 \cdot 1  \; \mapsto \; 1 
  = \frac{2}{2} &\\
  {\phantom{\funct{f} :}}& \; 2 \; \mapsto \; 4 = 2 \cdot 2  \; \mapsto \; 2 
  = \frac{4}{2} &
\end{flalign*}
\begin{flalign*}
  {\phantom{\funct{f} :}}& \; 3 \; \mapsto \; 6 = 2 \cdot 3  \; \mapsto \; 3 
  = \frac{6}{2} &\\
  {\phantom{\funct{f} :}}& \; 4 \; \mapsto \; 8 = 2 \cdot 4  \; \mapsto \; 4 
  = \frac{8}{2} &
\end{flalign*}
\begin{flalign*}
  {\phantom{\funct{f} :}}& \; 5 \; \mapsto \;10 = 2 \cdot 5  \; \mapsto \; 5 
  = \frac{10}{2} &\\
  {\phantom{\funct{f} :}}& \; 6 \; \mapsto \;12 = 2 \cdot 6  \; \mapsto \; 6 
  = \frac{12}{2} &
\end{flalign*}
\begin{flalign*}
  {\phantom{\funct{f} :}}& \; 7 \; \mapsto \;12 = 2 \cdot 7  \; \mapsto \; 7 
  = \frac{14}{2} &\\
  {\phantom{\funct{f} :}}& \; 8 \; \mapsto \;12 = 2 \cdot 8  \; \mapsto \; 8 
  = \frac{16}{2} &
\end{flalign*}
\begin{flalign*}
  {\phantom{\funct{f} :}}& \; 9 \; \mapsto \;12 = 2 \cdot 9  \; \mapsto \; 9 
  = \frac{18}{2} &\\
  {\phantom{\funct{f} :}}& ... &\\
  {\phantom{\funct{f} :}}& \; n \; \mapsto \;n_{par} = 2 \cdot n  \; \mapsto 
  \; n = \frac{2\cdot n}{2} = \frac{n_{par}}{2} &
\end{flalign*}


	Como resumen $\funct{f}\circ\functinv{f} = \mathds{1}_{\setN}$, esto es, 
la identidad en $\setN$, formalmente, $n \in \setN \Rightarrow 
\funct{f}\circ\functinv{f}\left(n\right) = n$.


	Con lo que acabamos de ver, queda demostrado que $\funct{f}$ es función 
que tiene inversa $\functinv{f}$, ambas son biyectivas y 
$\funct{f}\circ\functinv{f} = \mathds{1}_{\setN} : \setN \fromto \setN$.


	Pero aún vamos a verificar la otra composición de aplicaciones, 
$\functinv{f}\circ\funct{f}$ y veremos que coincide con 
$\mathds{1}_{\textsf{Par}}$, la función unidad en el conjunto de los pares:



\begin{flalign*}
  \functinv{f}\circ\funct{f} :\; & \mathsf{Par} \quad \fromto \quad \setN 
  \quad \fromto \quad \mathsf{Par}& \\
  {\phantom{\functinv{f}\circ\funct{f} :\; }}& \; 0 \; \mapsto \; 0 = 
  \frac{0}{2}  \; \mapsto \; 0 = 2 \cdot 0 &
\end{flalign*}
\begin{flalign*}
  {\phantom{\functinv{f}\circ\funct{f} :\; }}& \; 2 \; \mapsto \; 1 = 
  \frac{2}{2}  \; \mapsto \; 2 = 2 \cdot 1 &\\
  {\phantom{\functinv{f}\circ\funct{f} :\; }}& \; 4 \; \mapsto \; 2 = 
  \frac{4}{2}  \; \mapsto \; 4 = 2 \cdot 2 &
\end{flalign*}
\begin{flalign*}
  {\phantom{\functinv{f}\circ\funct{f} :\; }}& \; 6 \; \mapsto \; 3 = 
  \frac{6}{2}  \; \mapsto \; 6 = 2 \cdot 3 &\\
  {\phantom{\functinv{f}\circ\funct{f} :\; }}& \; 8 \; \mapsto \; 4 = 
  \frac{8}{2}  \; \mapsto \; 8 = 2 \cdot 4 &
\end{flalign*}
\begin{flalign*}
  {\phantom{\functinv{f}\circ\funct{f} :\; }}& \; 10 \; \mapsto \; 5 = 
  \frac{10}{2}  \; \mapsto \; 10 = 2 \cdot 5 &\\
  {\phantom{\functinv{f}\circ\funct{f} :\; }}& \; 12 \; \mapsto \; 6 = 
  \frac{12}{2}  \; \mapsto \; 12 = 2 \cdot 6 &
\end{flalign*}
\begin{flalign*}
  {\phantom{\functinv{f}\circ\funct{f} :\; }}& \; 14 \; \mapsto \; 7 = 
  \frac{14}{2}  \; \mapsto \; 14 = 2 \cdot 7 &\\
  {\phantom{\functinv{f}\circ\funct{f} :\; }}& \; 16 \; \mapsto \; 8 = 
  \frac{16}{2}  \; \mapsto \; 16 = 2 \cdot 8 &
\end{flalign*}
\begin{flalign*}
  {\phantom{\functinv{f}\circ\funct{f} :\; }}& \; 18 \; \mapsto \; 9 = 
  \frac{18}{2}  \; \mapsto \; 18 = 2 \cdot 9 &
\end{flalign*}
\begin{flalign*}
  {\phantom{\functinv{f}\circ\funct{f} :\; }}& ... &\\
  {\phantom{\functinv{f}\circ\funct{f} :\; }}& \; n_{par} \; \mapsto \;n = 
  \frac{n_{par}}{2} = \frac{2 \cdot n}{2} = n  \; \mapsto \; n_{par} = 2\cdot 
  n&
\end{flalign*}


	Ahora ya no debe quedar ninguna duda: los conjuntos $\setN$ y 
$\textsf{Par}$ tienen el mismo cardinal, el nombre que se le suele dar es 
$\aleph_0$. El subíndice $0$ viene del hecho que es el cardinal más pequeño 
que hay entre los cardinales infinitos.



	Por extraño que parezca, existen muchísimos subconjuntos de $\setN$, 
cualquiera que sea infinito, que tiene su mismo cardinal: $\aleph_0 \defin 
\card{\setN}$.



	Queda por ver si realmente existen conjuntos más grandes que $\setN$. Por 
ejemplo, podemos compararlo con un conjunto que sabemos que tiene más 
elementos, como por ejemplo $\setZ$ y $\setQ$.



	Por ejemplo ¿seremos capaces de encontrar una función biyectiva que vaya 
de $\setZ$ a $\setN$? Por ejemplo: ordenemos los números enteros (elementos 
del conjunto origen $\setZ$) de la siguiente forma:


$$
\begin{matrix}
  0 & 1 & -1 & 2 & -2 & 3 \\
  -3 & 4 & -4 & 5 & -5 & 6 \\
  -6 & 7 & -7 & 8 & -8 & 9 \\
  \ldots & \ldots & \ldots & \ldots & \ldots & \ldots \\
  -\iota & \iota+1 & -\iota+1 & \iota+2 & -\iota+2 & \iota+3 \\
  \ldots & \ldots & \ldots & \ldots & \ldots & \ldots \\
\end{matrix}
$$


	Ahora ordenamos los del conjunto final $\setN$ en su forma habitual;


$$
\begin{matrix}
  0 & 1 & 2 & 3 & 4 & 5 \\
  6 & 7 & 8 & 9 & -5 & 10 \\
  11 & 12 & 13 & 14 & 15 & 16 \\
  \ldots & \ldots & \ldots & \ldots & \ldots & \ldots \\
  -\iota & \iota+1 & -\iota+1 & \iota+2 & -\iota+2 & \iota+3 \\
  \ldots & \ldots & \ldots & \ldots & \ldots & \ldots \\
\end{matrix}
$$



Ahora tomaremos las siguientes reglas a cualquier $z \in \setZ$:



$$
\functat{g}{z} =
\left\{
  \begin{array}{ll}
    -2 \cdot z , & \hbox{si $z \leq 0$;} \texttt{ Salen números pares para 
    los no positivos.}\\
    2 \cdot z - 1, & \hbox{si $z \lneq 0$.} \texttt{ Salen numeros impares 
    para los negativos.}
  \end{array}
\right.
$$


Si cogéis los elementos del conjuntos origen, los enteros, en el orden 
explicado antes, y le aplicáis la regla dada por $\funct{g}$ de la última 
fórmula, obtenéis los números naturales en el orden dado anteriormente. 
$\funct{g} : \setZ \fromto \setN$ \textbf{\emph{es función biyectiva}}. 
Además, $\functinv{g}$ existe y viene dada por:


$$
\functinvat{g}{n} =
\left\{
  \begin{array}{ll}
    -\frac{n}{2} , & \hbox{si $\exists k \in \setN \quad n =  2 \cdot k $;} 
    \texttt{ No positivos para los pares.}\\
    \frac{n+1}{2}, & \hbox{si $\exists k \in \setN \quad n =  2 \cdot k-1 $.} 
    \texttt{ Negativos para los impares.}
  \end{array}
\right.
$$


	Ahora vamos a componer $\funct{g}\circ\functinvat{g}{n} = 
\functat{g}{\functinvat{g}{n}}$. Desarrollemos en dos pasos y por casos la 
nueva función, dónde los casos son que sea $n$ par o impar.


Primer paso (teniendo en cuanta que $n$ es par o no):
$$
\functat{g}{\functinvat{g}{n}} =
\left\{
  \begin{array}{ll}
    \functat{g}{-\frac{n}{2}} , & n \in \textsf{Par}\text{;}\\
    \functat{g}{\frac{n+1}{2}}, & n \in \textsf{Impar}\text{.}
  \end{array}
\right.
$$


Segundo paso (teniendo en cuenta si el resultado anterior es negativo o 
positivo):
$$
\funct{g}\circ\functinv{g}\left({n}\right) =
\left\{
  \begin{array}{ll}
    -2\cdot{-\frac{n}{2}} , & n \in \textsf{Par}\text{;}\\
    2\cdot{\frac{n+1}{2}}-1, & n \in \textsf{Impar}\text{.}
  \end{array}
\right.
$$


Reuniéndolo todo. Conclusión:
$$
\funct{g}\circ\functinv{g}\left({n}\right) =
\left\{
  \begin{array}{ll}
    n=2\cdot{\frac{n}{2}}=-2\cdot{-\frac{n}{2}} , & n \in 
    \textsf{Par}\text{;}\\
    n=(n+1)-1=2\cdot{\frac{n+1}{2}}-1, & n \in \textsf{Impar}\text{.}
  \end{array}
\right.
$$
$$
{{\mathds{1}}_{\setN}{\left({n}\right)}=
\funct{g}\circ\functinv{g}}{\left({n}\right)} = n \quad \forall n \in \setN
$$


	También le podemos dar la vuelta. Queremos ver cómo es la composición 
$\functinv{g}\circ\functat{g}{z} = \functinvat{g}{\functat{g}{z}}$ para los 
elementos $z \in \setZ$.  Volvemos a desarrollar en dos pasos y por casos la 
nueva función, ahora los pasos son inversos a los anteriores, dónde los casos 
ahora son que sea $z$ sea positivo ó no positivo.




Primer paso (teniendo en cuenta que $z$ es positivo o no):
$$
\functinvat{g}{\functat{g}{z}} =
\left\{
  \begin{array}{ll}
    \functinvat{g}{-2 \cdot z} , & z \leq 0 \text{;}\;\;{ n \defin -2\cdot z} 
    \texttt{ y } {n \in \textsf{Par} \subset \setN}\\
    \functinvat{g}{2 \cdot z - 1}, & z > 0 \text{.}\;\;{ n \defin 2\cdot z - 
    1} \texttt{ y } {n \in \textsf{Impar} \subset \setN}
  \end{array}
\right.
$$


Segundo paso (teniendo en cuenta si el resultado anterior $n$ es par o 
impar):
$$
\functinv{g}\circ\funct{g}\left({z}\right) =
\left\{
  \begin{array}{ll}
    z = \frac{-2 \cdot z}{2} , & z \leq 0 \text{;}\;\;{ n \defin -2\cdot z} 
    \texttt{ y } {n \in \textsf{Par} \subset \setN}\\
    z = \frac{{2 \cdot z - 1} + 1}{2} , & z > 0 \text{.}\;\;{ n \defin 2\cdot 
    z - 1} \texttt{ y } {n \in \textsf{Impar} \subset \setN}
  \end{array}
\right.
$$


Reuniéndolo todo:
$$
\functinv{g}\circ\funct{g}\left({z}\right) =
\left\{
  \begin{array}{ll}
    z , & z \leq 0 \text{;}\;\;{z \in \setZ^{-}\cup\set{0} \subset \setZ}\\
    z , & z > 0 \text{;}\text{.}\;\;{z \in \setZ^{+} \subset \setZ}
  \end{array}
\right.
$$
$$
{{\mathds{1}}_{\setZ}{\left({z}\right)}=
\functinv{g}\circ\funct{g}}{\left({z}\right)} = z\quad \forall z \in \setZ
$$


	Ahora hemos demostrado que $\card{\setZ} = \card{\setN} = 
\card{\textsf{Par}} = \card{\textsf{Impar}}$ o dicho con la relación de 
equipotencia, $\setZ \simeq \setN \simeq \textsf{Par} \simeq \textsf{Impar}$. 
Extraño, realmente extraño. Más adelante indagaremos en qué es lo que nos 
parece tan extraño de todos esto.


	Pero por ahora seguiremos buscando un conjunto que sea realmente más 
grande que $\setN$. Vamos a probar suerte con $\setQ^{+,0} \defin \setQ^+ 
\cup \set{0}$, esto es con los números racionales que son positivos o cero.


	Veamos algunas propiedades relevantes. Entre dos números consecutivos 
naturales no hay ningún otro número. ¿Porqué nos puede importar eso? Porque 
entre esos mismos dos números llevados a $\setQ$ hay una cantidad infinita de 
números. Y como el infinito más pequeño es el de los conjuntos de cardinal 
$\aleph_0$, entre cualquiera dos números consecutivos en $\setN$ encontramos 
en $\setQ$ un conjunto de cardinal $\aleph_0$. Pero me quedo corto. 
Concretemos, entre los números $0$ y $1$ tenemos que 
$\card{\setprop{x\in\setN}{1 > x > 0}}=\card{\emptyset}=0$. Pero $\forall k 
\in \setN \;\; \card{\setprop{x\in\setQ}{1 > x > 0}} > k$. O sea, que el 
conjunto que llamaremos ${\left]0,1\right[}_{\setQ} \defin 
\setprop{x\in\setQ}{1 > x > 0}$, es al menos tan grande como $\setN$. Pero 
pongamos ahora ${\left]0,0.1\right[}_{\setQ} \defin \setprop{x\in\setQ}{0.1 > 
x > 0}$. Podemos encontrar la siguiente serie dentro de ese conjunto 
$$\textsf{Serie}\defin\set{0.1-\frac{0.1}{2},0.1-\frac{2\cdot 
0.1}{3},0.1-\frac{3\cdot 0.1}{4},\ldots, 0.1-\frac{n\cdot 0.1}{n+1},\ldots}$$ 
que podemos ver que es tan grande como $\setN$ ya que los referenciados $n$ 
son todos los que hay en $\setN\setminus{0}$. Cojamos ahora dos elementos 
sucesivos de $\textsf{Serie}$ como $n_1 \defin 0.1-\frac{3\cdot 0.1}{4} = 
\frac{1}{40}$ y $n_2 \defin 0.1-\frac{2\cdot 0.1}{3} = \frac{2}{30} =  
\frac{1}{15}$. Ahora solo nos interesa que $n_1 < n_2$ y que podemos realizar 
una serie como $$\textsf{Serie}_2 \defin \set{\frac{n_1 + n_2}{2}, \frac{n_1 
+ 2 \cdot n_2}{3},\ldots,\frac{n_1 + \iota \cdot n_2}{\iota+1},\ldots}$$


	Es importante darse cuenta que $n_1$ y $n_2$ podrían haber sido dos 
números cualesquiera de $\setQ$ y siempre tendríamos que el conjunto contiene 
infinitos números como en $\textsf{Serie}$ y $\textsf{Serie}_2$. Cualquier 
intervalo entre dos números en $\setQ$ contiene una serie de números 
infinita.


	$\setQ$ parece un buen candidato para ser más grande que $\setN$, que 
hasta ahora ha permanecido imbatible.


\newcommand{\ordpair}[2]{\left({{#1},{#2}}\right)}


	Pero nos vamos a servir de un truco: todo número fraccionario, racional, 
se constituye de numerador y denominador. Así que usando la función $ 
\funct{\iota}: \setQ^{0,+} \fromto \setN^{0}\times\setN^{0} :: \frac{n}{m} 
\mapsto \ordpair{n}{m}$, que nos asegura que toda representación de un número 
racional positivo, incluyendo el cero, tiene una imagen distinta (propia) 
como pareja de números naturales (incluyendo el cero) ordenada. Esto es, 
$\setQ^{0,+} \leqpot \setN^{0}\times\setN^{0}$. Por lo tanto si conseguimos 
ver que existe alguna relación de cardinalidad entre $\setN^{0}$ y 
$\setN^{0}\times\setN^{0}$, quizás podamos ver qué es lo que pasa entre 
$\setN^{0}$ y $\setQ^{0,+}$ en términos de cardinalidad, de ver cuál es más
grande o si son igual de grandes.


	Veamos la siguiente construcción, dónde las cabeceras de columnas están 
etiquetadas con $f_n$ dónde ese $n$ corresponde con el primer componente del 
par, los pares de su columna, y en las filas, la etiqueta $f_m$, dónde $m$ 
corresponde al segundo componente de los pares ordenados de esa fila:


\newcommand{\ordpairlist}[1]{\ordpair{0}{#1} & \ordpair{1}{#1} &
                             \ordpair{2}{#1} & \ordpair{3}{#1} &
                             \ordpair{4}{#1} & \ordpair{5}{#1} &
                             \ordpair{6}{#1} & \ordpair{7}{#1} }
$$
\begin{matrix}
  {\setN}^{0}\times{\setN}^{0} & c_0 & c_1 & c_2 & c_3 & c_4 & c_5 & c_6 &
  c_7 & \ldots \\
  f_0 & \ordpairlist{0} & \ldots\\
  f_1 & \ordpairlist{1} & \ldots\\
  f_2 & \ordpairlist{2} & \ldots\\
  f_3 & \ordpairlist{3} & \ldots\\
  f_4 & \ordpairlist{4} & \ldots\\
  f_5 & \ordpairlist{5} & \ldots\\
  f_6 & \ordpairlist{6} & \ldots\\
  f_7 & \ordpairlist{7} & \ldots\\
  \vdots & \vdots & \vdots & \vdots & \vdots & \vdots & \vdots & \vdots & 
  \vdots & \ddots \\
\end{matrix}
$$


	Bien, si te das cuenta, en el esquema anterior, vemos que aparecen todas 
las posibles fracciones (no negativas, esto es , positivas o cero) pero 
muchas aparecen muchas veces, por ejemplo, la columna $c_0$ representa toda 
ella el número racional $0$ al ser el numerador $0$ en todas esas fracciones. 
Incluso, la fila $0$ no tiene ningún sentido como representación de 
racionales, pues el denominador no puede ser nunca $0$.
El siguiente ejemplo sería la diagonal que parte de $f_1 c_1$ y todas las 
casillas dónde el número subíndice de fila y columna coincidan: representan 
al racional $1$, al ser iguales numerador y denominador. Y podéis ver muchos 
otros ejemplos. Por ejemplo, la fila $f_1$ representa ella sola todos los 
números naturales incluyendo el $0$. El conjunto $\setQ^{0,+}$ está dentro de 
este conjunto representado en el esquema, conjunto que podemos denominar 
$\setN^0\times\setN^0$. A ésta, la operación 
representada por $\times$, la llamamos producto cartesiano de dos conjuntos.


	Como el producto cartesiano es una operación con la que nos hemos 
encontrado y nos encontraremos con frecuencia, la vamos a detallar un poco. 
Lo importante: forma un conjunto a partir de los dos conjuntos que 
intervienen en el producto cartesiano. También fundamental: los elementos del 
conjunto resultante son \textbf{\emph{pares ordenados}}, esto es, parejas de 
elementos. El primer elemento del par pertenece al primer conjunto del 
producto cartesiano, y el segundo elemento de par pertenece al segundo 
conjunto del producto. Además, ¿cómo caracterizamos un par ordenado? Es 
fácil. $\left({a,b}\right)=\left({c,d}\right)$ si y solo si $a = c$ y $b = 
d$. Con lo dicho bastaría, pero estamos tratando siempre con conjuntos y un 
par ordenado no parece un conjunto. Sin embargo hay un montón de formas de 
conseguir hacer pares ordenados con conjuntos, el más común es el siguiente:


$$
\forall x \in \sX \quad \forall y \in \sY \qquad \ordpair{x}{y} \in 
\sX\times\sY
$$
es un conjunto tal que
$$
\ordpair{x}{y} \defin \set{\set{x},\set{x,y}}
$$


	Si llamamos ${\funct{p}_1} : {\sX\times\sY} \fromto {\sX}$ a la función 
que nos da el primer elemento del par (este tipo de funciones se llaman 
proyecciones), podríamos concretarla así: 
${\funct{p}_1}{\left({\left({x,y}\right)}\right)}$ sería igual a $x$ que es 
el elemento único de 
$$\bigcup{\left({x,y}\right)} = \bigcup{\set{\set{x},\set{x,y}}} = 
\set{x}\cap\set{x,y} = \set{x}$$
De otro lado ${\funct{p}_2} : {\sX\times\sY} \fromto {\sY}$ lo podríamos 
formalizar como el único elemento de:
$$\set{x}\triangle\set{x,y}\defin\left({\set{x}}\setminus{\set{x,y}}\right)\cup
\left({\set{x,y}}\setminus{\set{x}}\right) = 
\set{}\cup\set{y} = \set{y}$$



	Pongamos que $\left({a,b}\right)=\left({c,d}\right)$ querría decir que 
$\set{\set{a},\set{a,b}} = \set{\set{c},\set{c,d}}$. Entonces $a = 
\funct{p}_1{\left({\left({a,b}\right)}\right)}=
\funct{p}_1{\left({\left({c,d}\right)}\right)}= c$ e igualmente $b = 
\funct{p}_2{\left({\left({a,b}\right)}\right)}=
\funct{p}_2{\left({\left({c,d}\right)}\right)}= d$.



	Bien, ya sabemos que es el esquema que hemos encontrado al representar 
todas las fracciones, como parejas de números, parejas con orden, dónde no da
lo mismo dónde se encuentren los números, en un conjunto producto cartesiano.
Podemos encontrar una función sobreyectiva de
$\setN^{0}\times\setN\fromto\setQ^{0,+}$, y esta es
$\funct{f}((n,m)) = \frac{p}{q}$ si $\frac{p}{q} = 
\text{Irreducible}({\frac{n}{m}})$. Esto indica que 
$\setN^{0}\times\setN\geqpot\setQ^{0,+}$.



	Ahora vamos a encontrar una forma de poner en una sola línea todos los 
elementos de $\setN^{0}\times\setN^{0}$. En el esquema anterior, vemos que la 
esquina superior izquierda $f_0,c_0$ no corresponde con ningún número
racional. De hecho, toda la fila $f_0$ no representa a ningún número 
racional. La $f_0,c_0$ representada por $\ordpair{0}{0}$ dónde la 
suma de sus proyecciones es $0$, no hay más elementos que su
suma sea $0$. Ahora podemos ver la los elementos de la tabla $f_0,c_1$ y
$f_1,c_0$ que son los pares de suma $1$. No hay más pares con esa suma. Ahora 
podemos encontrar aquellas parejas que su suma sea $2$. Son pocas 
$\ordpair{0}{2},\ordpair{1}{1},\ordpair{2}{1}$. Ya 
no hay más de suma $2$. Ahora las de suma $3$,
$\ordpair{0}{3},\ordpair{1}{2},\ordpair{2}{1},\ordpair{3}{0}$. Las
de suma $4$ son
$\ordpair{0}{4},\ordpair{1}{3},\ordpair{2}{2},\ordpair{3}{1},\ordpair{4}{0}$.


	Y así podemos seguir...


	Para suma $0$ tenemos $1$ solo elemento. Será el primer elemento de
nuestra serie. $\textsf{s}_0 = \ordpair{0}{0}$. Para suma $1$ hay $2$
elementos para la serie $\textsf{s}_1 = \ordpair{0}{1}$ y  $\textsf{s}_2 =
\ordpair{1}{0}$.
Para suma $2$ tenemos $3$ elementos, $\textsf{s}_3 = \ordpair{0}{2}$,
$\textsf{s}_4 = \ordpair{1}{1}$ y $\textsf{s}_5 = \ordpair{2}{0}$. Para suma
$3$ tenemos $4$ pares ordenados cuyas proyecciones (componentes, partes) 
sumen $3$, $\textsf{s}_6 = \ordpair{0}{3}$,  $\textsf{s}_7 = \ordpair{1}{2}$,
$\textsf{s}_8 = \ordpair{2}{1}$,  $\textsf{s}_9 = \ordpair{3}{0}$. Como
último ejemplo, pondremos los siguientes, los de suma $4$. Estos son
$\textsf{s}_{10} = \ordpair{0}{4}$,  $\textsf{s}_{11} = \ordpair{1}{3}$,
$\textsf{s}_{12} = \ordpair{2}{2}$,  $\textsf{s}_{13} = \ordpair{3}{1}$ y por
último $\textsf{s}_{14} = \ordpair{4}{0}$.


\begin{multline*}
  \quad \qquad \qquad \qquad \qquad \qquad \qquad \qquad \textsf{s}_{0} =
  \ordpair{0}{0} \; \\
  \textsf{s}_{1} = \ordpair{0}{1} \;  \textsf{s}_{2} = \ordpair{1}{0}  \\
  \textsf{s}_{3} = \ordpair{0}{2} \;  \textsf{s}_{4} = \ordpair{1}{1}  \;
  \textsf{s}_{5} = \ordpair{2}{0} \;\\
  \textsf{s}_{6} = \ordpair{0}{3} \;  \textsf{s}_{7} = \ordpair{1}{2}  \;
  \textsf{s}_{8} = \ordpair{2}{1} \;  \textsf{s}_{9} = \ordpair{3}{0}  \;
  \\
  \textsf{s}_{10} = \ordpair{0}{4} \;  \textsf{s}_{11} = \ordpair{1}{3}  \;
  \textsf{s}_{12} = \ordpair{2}{2} \;  \textsf{s}_{13} = \ordpair{3}{1}  \;
  \textsf{s}_{14} = \ordpair{4}{0} \;
  \\
  \textsf{s}_{15} = \ordpair{0}{5} \;  \textsf{s}_{16} = \ordpair{1}{4}  \;
  \textsf{s}_{17} = \ordpair{2}{3} \;  \textsf{s}_{18} = \ordpair{3}{2}  \;
  \textsf{s}_{19} = \ordpair{4}{1} \;  \textsf{s}_{20} = \ordpair{5}{0}  \;
  \\
  \ldots \;  \ldots  \;  \ldots \;  \ldots  \;  \ldots \;  \ldots \; \ldots
  \; \ldots \;  \ldots  \;  \ldots \;  \ldots  \;  \ldots \;  \ldots \;
  \ldots \; \ldots \;  \ldots  \;  \ldots \;  \ldots  \;  \ldots \;  \ldots
  \; \ldots \;
\end{multline*}


	La primera línea corresponde a los pares ordenados de suma $0$. La
llamaremos grupo $g_0$. Hay $1$ par. $\numeral g_0 = 1$.


	La segunda línea corresponde a los pares ordenados de suma $1$. La
llamaremos grupo $g_1$. Hay $2$ pares. $\numeral g_1 = 2$.


	La tercera línea corresponde a los pares ordenados de suma $2$. La
llamaremos grupo $g_2$. Hay $3$ pares. $\numeral g_2 = 3$.


	La cuarta línea corresponde a los pares ordenados de suma $3$. La
llamaremos grupo $g_3$. Hay $4$ pares. $\numeral g_3 = 4$.


	La quinta línea corresponde a los pares ordenados de suma $4$. La
llamaremos grupo $g_4$. Hay $5$ pares. $\numeral g_4 = 5$.


	La sexta línea corresponde a los pares ordenados de suma $5$. La
llamaremos grupo $g_5$. Hay $6$ pares. $\numeral g_5 = 6$.


	La enésima línea corresponde a los pares ordenados de suma $n-1$. La
llamaremos grupo $g_{n-1}$. Hay $n$ pares. $\numeral g_{n-1} = n$.


	¿Cuántos pares ordenados hay desde el principio hasta la sexta línea
incluida? Habrá $1+2+3+4+5+6 = \numeral g_0 + \numeral g_1 + \numeral g_2 +
\numeral g_3 + \numeral g_4 + \numeral g_5$ pares ordenados, esto es, $6
\cdot 7 / 2 \; = \; 21$ pares ordenados (para incluir el grupo de suma
$\numeral g_5 = 6$), según una conocida fórmula de las sumas de los primeros
$\nu$ números: $1+2+\ldots+\nu \; = \; \nu \cdot (\nu + 1) / 2$. Aplicada a
los grupos de misma suma, $\numeral g_0 = 1$ par ordenado, para el índice
$0$, para los dos primeros grupos $g_0 , g_1$ tenemos $2\cdot 3/ 2 = 3$ pares
ordenados en total, comenzando por el índice $0$ del $g_0$, y los índices $1$
y $2$ para el grupo $g_1$. Para los tres primeros grupos $g_0, g_1, g_2$
tenemos $3\cdot 4 / 2 = 6$, dónde el grupo $g_2$ va desde el índice $3$ hasta
el $5$. Y, en general, para los $\nu$ primeros grupos, tenemos $(\nu + 1)
\cdot (\nu + 2)/2$ elementos de la serie en total, siendo el grupo $g_{\nu}$
el que va del índice $\nu \cdot \left(\nu + 1\right) / 2$ hasta el $\nu \cdot
\left(\nu + 3\right) / 2$.


	Conociendo esto podemos calcular para cada $g_{\nu}$ cuál es el índice
mayor y el menor del grupo. El índice menor del grupo será: $\nu \cdot (\nu +
1) / 2$ y el mayor $(\nu + 3) \cdot \nu / 2$. Así un índice $\iota$
cualquiera de $\setN$ estará en un grupo $g_{\nu}$, de forma que


$$
\nu \cdot (\nu + 1) / 2 \leq \iota \leq \nu \cdot (\nu + 3) / 2
$$,


	y de esta desigualdad, multiplicando por $2$,


$$
{\nu}^2 + {\nu} \leq 2 \cdot \iota \leq {\nu}^2 + 3 \cdot \nu
$$. 


	De aquí saldrán dos ecuaciones:


$$
{\nu}^2 + {\nu} - 2 \cdot \iota = 0
$$ y 
$$
{\nu}^2 + 3 \cdot {\nu} - 2 \cdot \iota = 0
$$.


	Que son sencillas ecuaciones cuadráticas, con soluciones que deben ser
positivas y enteras:


$$
\functat{{\nu}_{\min}}{\iota} =
\left\lfloor\frac{\sqrt{1+8\cdot\iota}-1}{2}\right\rfloor
$$ y 
$$
\functat{{\nu}_{\max}}{\iota} =
\left\lceil\frac{\sqrt{9+8\cdot\iota}-3}{2}\right\rceil
$$.


	Hagamos unos pocos cálculos:


Cálculos para $\iota = 0$:


$\functat{{\nu}_{\min}}{0} =
\left\lfloor\frac{\sqrt{1+8\cdot 0}-1}{2}\right\rfloor =
\left\lfloor\frac{\sqrt{1}-1}{2}\right\rfloor =
\left\lfloor\frac{0}{2}\right\rfloor =
\left\lfloor{0}\right\rfloor = 0
$


$\functat{{\nu}_{\max}}{0} =
\left\lceil\frac{\sqrt{9+8\cdot 0}-3}{2}\right\rceil =
\left\lceil\frac{0}{2}\right\rceil =
\left\lceil{0}\right\rceil = 0
$


$\functat{{\nu}_{\min}}{0} = 0$ y $\functat{{\nu}_{\max}}{0} = 0$


Cálculos para $\iota = 1$:


$\functat{{\nu}_{\min}}{1} =
\left\lfloor\frac{\sqrt{1+8\cdot 1}-1}{2}\right\rfloor =
\left\lfloor\frac{2}{2}\right\rfloor =
\left\lfloor{1}\right\rfloor = 1
$


$\functat{{\nu}_{\max}}{1} =
\left\lceil\frac{\sqrt{9+8\cdot 1}-3}{2}\right\rceil =
\left\lceil\frac{\sqrt{17}-3}{2}\right\rceil =
\left\lceil\frac{1.123...}{2}\right\rceil =
\left\lceil{0.56...}\right\rceil = 1
$


$\functat{{\nu}_{\min}}{1} = 1$ y $\functat{{\nu}_{\max}}{1} = 1$


Cálculos para $\iota = 2$:


$\functat{{\nu}_{\min}}{2} =
\left\lfloor\frac{\sqrt{1+8\cdot 2}-1}{2}\right\rfloor =
\left\lfloor\frac{\sqrt{17}-1}{2}\right\rfloor =
\left\lfloor\frac{2.12310...}{2}\right\rfloor =
\left\lfloor{1.06155...}\right\rfloor = 1
$


$\functat{{\nu}_{\max}}{2} =
\left\lceil\frac{\sqrt{9+8\cdot 2}-3}{2}\right\rceil =
\left\lceil\frac{\sqrt{25}-3}{2}\right\rceil =
\left\lceil{1}\right\rceil = 1
$


$\functat{{\nu}_{\min}}{2} = 1$ y $\functat{{\nu}_{\max}}{2} = 1$


Cálculos para $\iota = 3$:


$\functat{{\nu}_{\min}}{3} =
\left\lfloor\frac{\sqrt{1+8\cdot 3}-1}{2}\right\rfloor =
\left\lfloor\frac{4}{2}\right\rfloor =
\left\lfloor{2}\right\rfloor = 2
$


$\functat{{\nu}_{\max}}{3} =
\left\lceil\frac{\sqrt{9+8\cdot 3}-3}{2}\right\rceil =
\left\lceil\frac{\sqrt{33}-3}{2}\right\rceil =
\left\lceil\frac{2.7445...}{2}\right\rceil =
\left\lceil{1.3722...}\right\rceil = 2
$


$\functat{{\nu}_{\min}}{3} = 2$ y $\functat{{\nu}_{\max}}{3} = 2$


Cálculos para $\iota = 4$:


$\functat{{\nu}_{\min}}{4} =
\left\lfloor\frac{\sqrt{1+8\cdot 4}-1}{2}\right\rfloor =
\left\lfloor\frac{\sqrt{33}-1}{2}\right\rfloor =
\left\lfloor\frac{4.7445...}{2}\right\rfloor =
\left\lfloor{2.372...}\right\rfloor = 2
$


$\functat{{\nu}_{\max}}{4} =
\left\lceil\frac{\sqrt{9+8\cdot 4}-3}{2}\right\rceil =
\left\lceil\frac{\sqrt{41}-3}{2}\right\rceil =
\left\lceil\frac{3.4031...}{2}\right\rceil =
\left\lceil{1.7015...}\right\rceil = 2
$


$\functat{{\nu}_{\min}}{4} = 2$ y $\functat{{\nu}_{\max}}{4} = 2$


Cálculos para $\iota = 5$:


$\functat{{\nu}_{\min}}{5} =
\left\lfloor\frac{\sqrt{1+8\cdot 5}-1}{2}\right\rfloor =
\left\lfloor\frac{\sqrt{41}-1}{2}\right\rfloor =
\left\lfloor\frac{5.403...}{2}\right\rfloor =
\left\lfloor{2.701...}\right\rfloor = 2
$


$\functat{{\nu}_{\max}}{5} =
\left\lceil\frac{\sqrt{9+8\cdot 5}-3}{2}\right\rceil =
\left\lceil\frac{\sqrt{49}-3}{2}\right\rceil =
\left\lceil{2}\right\rceil = 2
$


$\functat{{\nu}_{\min}}{5} = 2$ y $\functat{{\nu}_{\max}}{5} = 2$


Cálculos para $\iota = 6$:


$\functat{{\nu}_{\min}}{6} =
\left\lfloor\frac{\sqrt{1+8\cdot 6}-1}{2}\right\rfloor =
\left\lfloor\frac{\sqrt{49}-1}{2}\right\rfloor =
\left\lfloor{3}\right\rfloor = 3
$


$\functat{{\nu}_{\max}}{6} =
\left\lceil\frac{\sqrt{9+8\cdot 5}-3}{2}\right\rceil =
\left\lceil\frac{\sqrt{57}-3}{2}\right\rceil =
\left\lceil\frac{4.5498...}{2}\right\rceil =
\left\lceil{2.2749...}\right\rceil = 3
$


$\functat{{\nu}_{\min}}{6} = 3$ y $\functat{{\nu}_{\max}}{6} = 3$


Cálculos para $\iota = 7$:


$\functat{{\nu}_{\min}}{7} =
\left\lfloor\frac{\sqrt{1+8\cdot 7}-1}{2}\right\rfloor =
\left\lfloor\frac{\sqrt{57}-1}{2}\right\rfloor =
\left\lfloor\frac{6.5498...}{2}\right\rfloor =
\left\lfloor{3.2749...}\right\rfloor = 3
$


$\functat{{\nu}_{\max}}{7} =
\left\lceil\frac{\sqrt{9+8\cdot 7}-3}{2}\right\rceil =
\left\lceil\frac{\sqrt{65}-3}{2}\right\rceil =
\left\lceil\frac{5.06225...}{2}\right\rceil =
\left\lceil{2.5311...}\right\rceil = 3
$


$\functat{{\nu}_{\min}}{7} = 3$ y $\functat{{\nu}_{\max}}{7} = 3$


Cálculos para $\iota = 8$:


$\functat{{\nu}_{\min}}{8} =
\left\lfloor\frac{\sqrt{1+8\cdot 8}-1}{2}\right\rfloor =
\left\lfloor\frac{\sqrt{65}-1}{2}\right\rfloor =
\left\lfloor\frac{7.0622...}{2}\right\rfloor =
\left\lfloor{3.5311...}\right\rfloor = 3
$


$\functat{{\nu}_{\max}}{8} =
\left\lceil\frac{\sqrt{9+8\cdot 8}-3}{2}\right\rceil =
\left\lceil\frac{\sqrt{73}-3}{2}\right\rceil =
\left\lceil\frac{5.544...}{2}\right\rceil =
\left\lceil{2.772...}\right\rceil = 3
$


$\functat{{\nu}_{\min}}{8} = 3$ y $\functat{{\nu}_{\max}}{8} = 3$


Cálculos para $\iota = 9$:


$\functat{{\nu}_{\min}}{9} =
\left\lfloor\frac{\sqrt{1+8 \cdot 9}-1}{2}\right\rfloor =
\left\lfloor\frac{\sqrt{73}-1}{2}\right\rfloor =
\left\lfloor\frac{7.0622...}{2}\right\rfloor =
\left\lfloor{3.5311...}\right\rfloor = 3
$


$\functat{{\nu}_{\max}}{9} =
\left\lceil\frac{\sqrt{9+8\cdot 9}-3}{2}\right\rceil =
\left\lceil\frac{\sqrt{81}-3}{2}\right\rceil =
\left\lceil{3}\right\rceil = 3
$


$\functat{{\nu}_{\min}}{9} = 3$ y $\functat{{\nu}_{\max}}{9} = 3$


Cálculos para $\iota = 10$:


$\functat{{\nu}_{\min}}{10} =
\left\lfloor\frac{\sqrt{1+8 \cdot 10}-1}{2}\right\rfloor =
\left\lfloor\frac{\sqrt{81}-1}{2}\right\rfloor =
\left\lfloor{4}\right\rfloor = 4
$


$\functat{{\nu}_{\max}}{10} =
\left\lceil\frac{\sqrt{9+8\cdot 10}-3}{2}\right\rceil =
\left\lceil\frac{\sqrt{89}-3}{2}\right\rceil =
\left\lceil\frac{6.433...}{2}\right\rceil =
\left\lceil{3.216...}\right\rceil = 4
$


$\functat{{\nu}_{\min}}{10} = 4$ y $\functat{{\nu}_{\max}}{10} = 4$


Cálculos para $\iota = 11$:


$\functat{{\nu}_{\min}}{11} =
\left\lfloor\frac{\sqrt{1+8 \cdot 11}-1}{2}\right\rfloor =
\left\lfloor\frac{\sqrt{89}-1}{2}\right\rfloor =
\left\lfloor\frac{8.433...}{2}\right\rfloor =
\left\lfloor{4.416...}\right\rfloor = 4
$


$\functat{{\nu}_{\max}}{11} =
\left\lceil\frac{\sqrt{9+8\cdot 11}-3}{2}\right\rceil =
\left\lceil\frac{\sqrt{97}-3}{2}\right\rceil =
\left\lceil\frac{6.8488...}{2}\right\rceil =
\left\lceil{3.4244...}\right\rceil = 4
$


$\functat{{\nu}_{\min}}{11} = 4$ y $\functat{{\nu}_{\max}}{11} = 4$


Cálculos para $\iota = 12$:


$\functat{{\nu}_{\min}}{12} =
\left\lfloor\frac{\sqrt{1+8 \cdot 12}-1}{2}\right\rfloor =
\left\lfloor\frac{\sqrt{97}-1}{2}\right\rfloor =
\left\lfloor\frac{8.8488...}{2}\right\rfloor =
\left\lfloor{4.4244...}\right\rfloor = 4
$


$\functat{{\nu}_{\max}}{12} =
\left\lceil\frac{\sqrt{9+8\cdot 12}-3}{2}\right\rceil =
\left\lceil\frac{\sqrt{105}-3}{2}\right\rceil =
\left\lceil\frac{7.24695...}{2}\right\rceil =
\left\lceil{3.623475...}\right\rceil = 4
$


$\functat{{\nu}_{\min}}{12} = 4$ y $\functat{{\nu}_{\max}}{12} = 4$


Cálculos para $\iota = 13$:


$\functat{{\nu}_{\min}}{13} =
\left\lfloor\frac{\sqrt{1+8 \cdot 13}-1}{2}\right\rfloor =
\left\lfloor\frac{\sqrt{105}-1}{2}\right\rfloor =
\left\lfloor\frac{9.24695...}{2}\right\rfloor =
\left\lfloor{4.6234...}\right\rfloor = 4
$


$\functat{{\nu}_{\max}}{13} =
\left\lceil\frac{\sqrt{9+8\cdot 13}-3}{2}\right\rceil =
\left\lceil\frac{\sqrt{113}-3}{2}\right\rceil =
\left\lceil\frac{7.630...}{2}\right\rceil =
\left\lceil{3.815...}\right\rceil = 4
$


$\functat{{\nu}_{\min}}{13} = 4$ y $\functat{{\nu}_{\max}}{13} = 4$


Cálculos para $\iota = 14$:


$\functat{{\nu}_{\min}}{14} =
\left\lfloor\frac{\sqrt{1+8 \cdot 14}-1}{2}\right\rfloor =
\left\lfloor\frac{\sqrt{113}-1}{2}\right\rfloor =
\left\lfloor\frac{9.6301...}{2}\right\rfloor =
\left\lfloor{4.815...}\right\rfloor = 4
$


$\functat{{\nu}_{\max}}{14} =
\left\lceil\frac{\sqrt{9+8\cdot 14}-3}{2}\right\rceil =
\left\lceil\frac{\sqrt{121}-3}{2}\right\rceil =
\left\lceil{4}\right\rceil = 4
$


$\functat{{\nu}_{\min}}{14} = 4$ y $\functat{{\nu}_{\max}}{14} = 4$


Cálculos para $\iota = 15$:


$\functat{{\nu}_{\min}}{15} =
\left\lfloor\frac{\sqrt{1+8 \cdot 15}-1}{2}\right\rfloor =
\left\lfloor\frac{\sqrt{121}-1}{2}\right\rfloor =
\left\lfloor{5}\right\rfloor = 5
$


$\functat{{\nu}_{\max}}{15} =
\left\lceil\frac{\sqrt{9+8\cdot 15}-3}{2}\right\rceil =
\left\lceil\frac{\sqrt{129}-3}{2}\right\rceil =
\left\lceil\frac{8.3578...}{2}\right\rceil =
\left\lceil{4.1789...}\right\rceil = 5
$


$\functat{{\nu}_{\min}}{15} = 5$ y $\functat{{\nu}_{\max}}{15} = 5$


	Para el siguiente cálculo daremos un salto, y calcularemos por último $\iota = 20$ y $\iota = 21$:


$\functat{{\nu}_{\min}}{20} =
\left\lfloor\frac{\sqrt{1+8 \cdot 20}-1}{2}\right\rfloor =
\left\lfloor\frac{\sqrt{161}-1}{2}\right\rfloor =
\left\lfloor\frac{11.688...}{2}\right\rfloor =
\left\lfloor{5.844...}\right\rfloor = 5
$


$\functat{{\nu}_{\max}}{20} =
\left\lceil\frac{\sqrt{9+8\cdot 20}-3}{2}\right\rceil =
\left\lceil\frac{\sqrt{169}-3}{2}\right\rceil =
\left\lceil{5}\right\rceil = 5
$


$\functat{{\nu}_{\min}}{20} = 5$ y $\functat{{\nu}_{\max}}{20} = 5$


$\functat{{\nu}_{\min}}{21} =
\left\lfloor\frac{\sqrt{1+8 \cdot 21}-1}{2}\right\rfloor =
\left\lfloor\frac{\sqrt{169}-1}{2}\right\rfloor =
\left\lfloor{6}\right\rfloor = 6
$


$\functat{{\nu}_{\max}}{21} =
\left\lceil\frac{\sqrt{9+8\cdot 21}-3}{2}\right\rceil =
\left\lceil\frac{\sqrt{177}-3}{2}\right\rceil =
\left\lceil\frac{10.3041...}{2}\right\rceil =
\left\lceil{5.152...}\right\rceil = 6
$


$\functat{{\nu}_{\min}}{21} = 6$ y $\functat{{\nu}_{\max}}{21} = 6$


	Después de tanto calcular, no sé si os habréis dado cuenta que no hemos
demostrado nada: hemos conseguido ver que en los primeros casos nuestras
fórmulas funcionan bien. Y aunque creemos que van a funcionar siempre no lo
sabemos.


	Para demostrarlo, tendremos que volver a las fórmulas originales e
intentar ver para qué índices $\iota \in \setN$ podemos calcular mediante las
fórmulas como permanecen constantes los índices.


$$
\functat{{\nu^\prime}_{\min}}{\iota} =
\frac{\sqrt{1+8\cdot\iota}-1}{2}
$$ y 
$$
\functat{{\nu^\prime}_{\max}}{\iota} =
\frac{\sqrt{9+8\cdot\iota}-3}{2} =
\frac{\sqrt{1+8\cdot\left(\iota + 1\right)}-3}{2}
$$.


	Igualando las dos fórmulas podemos conseguir ver qué valores naturales
cumplen la igualdad:


$$
\frac{\sqrt{1+8\cdot\iota}-1}{2} =
\frac{\sqrt{1+8\cdot\left(\iota + 1\right)}-3}{2}
$$
$$
\sqrt{1+8\cdot\iota}-1 =
\sqrt{1+8\cdot\left(\iota + 1\right)}-3
$$
$$
\sqrt{1+8\cdot\iota} =
\sqrt{1+8\cdot\left(\iota + 1\right)}-2
$$
$$
{1+8\cdot\iota} = {1+8\cdot\left(\iota + 1\right)} + 4 - 4 \cdot
\sqrt{1+8\cdot\left(\iota + 1\right)}
$$
$$
0 = 12 - 4 \cdot \sqrt{1+8\cdot\left(\iota + 1\right)}
$$
$$
3 = \sqrt{1+8\cdot\left(\iota + 1\right)}
$$
$$
9 = {1+8\cdot\left(\iota + 1\right)}
$$
$$
\iota = 0
$$


	Como era de esperar, ya que exigimos la igualdad en $\setR$. Así que
intentaremos acotar la diferencia cuanto dos índices distintos toman
diferencias menores que $1$ y mayores que $0$:


$$
\imath > \jmath
$$


$$
0 < 
    \frac{\sqrt{1+8\cdot\imath}-1}{2}
    -
    \frac{\sqrt{1+8\cdot\jmath}-1}{2}
< 1
$$


	Multiplicando por $2$, cancelando los términos iguales $-1$ en ambos lados de la igualdad:


$$
\imath > \jmath
$$


$$
0 < \sqrt{1+8\cdot\imath}-\sqrt{1+8\cdot\jmath} < 2
$$


	Elevando al cuadrado todos los términos:


$$
0 <
1 + 8 \cdot \imath + 1 + 8 \cdot \jmath -
2 \cdot \sqrt{ 1 + 8 \cdot \imath } \cdot \sqrt{ 1 + 8 \cdot \jmath }
< 4
$$
$$
0 <
1 + 4 \cdot \left(\imath + \jmath\right) -
\sqrt{ 1 + 8 \cdot \imath } \cdot \sqrt{ 1 + 8 \cdot \jmath }
< 2
$$
$$
-1 <
4 \cdot \left(\imath + \jmath\right) -
\sqrt{ 1 + 8 \cdot \imath } \cdot \sqrt{ 1 + 8 \cdot \jmath }
< 1
$$














	Lo más interesante es que a partir del índice $\iota$ de $\setN^{0}$, el
elemento de la serie $\textsf{s}{\left({\iota}\right)}$ podemos calcular su
grupo $g_{\nu}$ y de ahí el mínimo y el máximo índice de su grupo.


	En general para suma $n$ tenemos $n+1$ elementos de la serie, supongamos
que nos encontramos en $\textsf{s}_{m-1}$. Entonces tendremos
$\textsf{s}_{m}=\ordpair{0}{m+1}$, $\textsf{s}_{m+1}=\ordpair{1}{m}$, ...,
$\textsf{s}_{m+n}=\ordpair{n}{m+1-n}$.



Así que nuestro único problema es encontrar la $l$ que anteriormente hemos
mencionado en la serie.

\newcommand {\mmin}[1]{l^{\min}_{#1}}
\newcommand {\mmax}[1]{l^{\max}_{#1}}
$$
\begin{matrix}
i = 1           & \mmin{1} = 1                & \mmax{1} = 1 \\
i = 2, 3        & \mmin{2} = \mmin{1} + 1 = 2 & \mmax{2} = \mmin{1} + 2 =  3
\\
i = 4, 5, 6     & \mmin{3} = \mmin{2} + 2 = 4 & \mmax{3} = \mmin{2} + 3 =  6
\\
i = 7, 8, 9, 10 & \mmin{4} = \mmin{3} + 3 = 7 & \mmax{4} = \mmin{3} + 4 = 10
\\
\ldots          & \ldots         & \ldots\\
i = \mmin{k-1}, ..., \mmin{k-1}+{k-1} & \mmin{k} = \mmin{k-1} + {k-1} &
\mmax{k} = \mmin{k-1} + k \\
\ldots & \ldots & \ldots
\end{matrix}
$$


	Ahora podemos tener una formulita para las $m_{\iota}$. $m_{\iota} =
m_{\iota - 1} + n_{\iota} = m_{\iota - 1} + \iota$. A este tipo de fórmulas
se les llama recursivas, y son muy importantes en programación y
computadores.


	La fórmula anterior, que parte de $m_1 = 2$ y $m_{\iota} = m_{\iota - 1}
+ \iota$, se puede poner como $m_1 = 1 + 1$, $m_2 = 2 + 2 = 4$, $m_3 = 4 + 3
= 7$, $m_4 = 7 + 4 = 11$, $m_{\kappa} = m_1 + n_2 + n_3 + \ldots +
n_{\kappa}$, $m_{\kappa} = 2 + 2 + 3 + \ldots + \kappa$,
$$
m_{\kappa} = 1 + (1 + 2 + 3 + 4 + \ldots + \kappa) = 1 + \frac{\kappa \cdot
\left({\kappa + 1}\right)}{2}
$$


	¿Para qué nos sirve esa fórmula? Supongamos que tenemos la fracción
$\frac{l}{n}$. Sabemos que esta en un segmento la $m = l + n$. Su índice será
$\frac{(l+n-1)(l+n)}{2}+l$. Así, $\frac{l}{n}$ ocupa el lugar $1 +
\frac{\kappa \cdot \left({\kappa + 1}\right)}{2}$ en la serie que estamos
formando.

$$
\functat{\iota}{\frac{l}{m}} = {1 + l + \frac{\left(l+m-1\right)
\cdot \left(l+m\right)}{2}}
$$











	% ----------------------------------------------------------------
	%\bibliographystyle{amsplain}
	%\bibliography{}
\end{document}
% ----------------------------------------------------------------
